\documentclass[11pt,a4paper]{article}
\usepackage[margin=23mm]{geometry}
\usepackage{fontspec}
\setmainfont{Latin Modern Roman}
\usepackage{amsmath,amssymb,amsthm,mathtools,booktabs,longtable,tabularx,array}
\usepackage[authoryear,round]{natbib}
\usepackage[colorlinks=true,linkcolor=blue,citecolor=blue,urlcolor=blue]{hyperref}
\usepackage{microtype}
\setlength{\emergencystretch}{3em}
\newtheorem{theorem}{Theorem}[section]
\newtheorem{proposition}[theorem]{Proposition}
\newtheorem{lemma}[theorem]{Lemma}
\newtheorem{corollary}[theorem]{Corollary}
\newtheorem{counterexample}[theorem]{Counterexample}
\newtheorem{example}[theorem]{Example}
\newtheorem{conjecture}[theorem]{Conjecture}
\theoremstyle{remark}
\newtheorem*{remark}{Remark}
\newcommand{\status}[1]{\textnormal{\textbf{#1}}}
\newcommand{\C}{\mathbb C}
\newcommand{\R}{\mathbb R}
\newcommand{\Z}{\mathbb Z}
\newcommand{\Q}{\mathbb Q}
\newcommand{\CC}{\mathbb C}
\newcommand{\RR}{\mathbb R}
\newcommand{\ZZ}{\mathbb Z}
\newcommand{\QQ}{\mathbb Q}
\newcommand{\PP}{\mathbb P}
\newcommand{\NN}{\mathbb N}
\newcommand{\TT}{\mathbb T}
\newcommand{\calD}{\mathcal D}
\newcommand{\B}{\mathrm B}
\DeclareMathOperator{\Exit}{Exit}
\DeclareMathOperator{\Map}{Map}
\DeclareMathOperator{\Fun}{Fun}
\DeclareMathOperator{\Hom}{Hom}
\DeclareMathOperator{\Aut}{Aut}
\DeclareMathOperator{\Stab}{Stab}
\DeclareMathOperator{\Shv}{Shv}
\DeclareMathOperator{\im}{im}
\DeclareMathOperator{\coker}{coker}
\DeclareMathOperator{\hofib}{hofib}
\DeclareMathOperator*{\hocolim}{hocolim}
\DeclareMathOperator{\codim}{codim}
\DeclareMathOperator{\diag}{diag}
\DeclareMathOperator{\Spec}{Spec}
\DeclareMathOperator{\rk}{rank}
\newcommand{\id}{\mathrm{id}}
\title{Exit categories, meridian lattices, and boundary reconstruction\\
\large A proof and counterexample audit of the six manuscript leads}
\author{AI-produced research draft: GPT-6 Astra, OpenAI\\
\small No claim of priority or independent publication}
\date{12 November 2026}
\begin{document}
\maketitle
\noindent\textbf{Conventions and scope, fixed before the statements.}
\begin{center}\small
\renewcommand{\arraystretch}{1.11}
\begin{tabularx}{\textwidth}{@{}p{0.21\textwidth}X@{}}
\toprule
Ambient objects&Complex manifolds, or smooth analytic Deligne--Mumford stacks presented by étale groupoids. Classifying spaces retain all discrete inertia. Proper finite-isotropy and algebraic hypotheses are imposed explicitly where orientation, Hodge theory, or factorization requires them.\\
Divisors and strata&Reduced normal-crossings divisors, including self-intersections and branch monodromy. Strata are connected open incidence strata; normalizations and finite branch-incidence covers retain the branches.\\
Exit convention&An exit moves from a deeper stratum to a more open one. All spaces and categories are derived; $B$ means groupoid completion or its classifying-space realization.\\
Groups and paths&Groups are discrete unless a completion is written. Basepoints and connecting paths are chosen; changing them conjugates the displayed maps. Lattices use positive complex meridians.\\
Goodness&Every stratum, including the open one, is aspherical, and each full normal-meridian lattice injects into the open group. This does not assert that the ambient space is aspherical.\\
Boundary and ends&$L=\partial X^{\log}$ is the resolved normal boundary. For proper pairs it represents the stable end. For nonproper pairs it need not represent every end of $U$.\\
Completions&The marking identifies $U=X\setminus D$ with the fixed open space. Algebraic markings and arbitrary holomorphic markings are distinguished. Analytic comparison is stated within a marked bimeromorphic class.\\
Blow-up centres&A smooth closed centre in $D$ is admissible when the centre and all local boundary branches are simultaneous coordinate subspaces. Disconnected centres are treated componentwise.\\
Curves and moduli&Stable types satisfy $2g-2+n>0$. Punctures are individually labelled; fixing boundary circles pointwise introduces the specified central twists. Ineffective stabilizers are retained.\\
Empty boundary&For $D=\varnothing$ there are no meridian tests and the normal boundary is empty. The exit category is the fundamental infinity-groupoid of the ambient space.\\
Arithmetic scope&The calculations are over $\C$ and their characteristic-zero comparison variants. No assertion about wild inertia or reduction at $2$ or $3$ is inferred from the cusp.\\
Reading policy&MPL package and scan numbers identify manuscript passages. Mathematical corrections are distinguished from transcription choices. A doubtful word is not declared paleographically certain because a theorem would be convenient.\\
\bottomrule
\end{tabularx}
\end{center}

\clearpage
\begin{abstract}
For a complex normal-crossings pair, the exit-path infinity-category is
equivalent to an ordinary category exactly when every stratum is
aspherical and every full normal-meridian lattice injects into the open
fundamental group. Adjacent tests alone do not suffice. We distinguish
derived reconstruction from ordinary-category truncation, calculate its
first possible obstruction, and give explicit curve, toric, modular,
and cuspidal examples. We analyze admissible blow-ups, good surface
completions, elementary fibrations, dual complexes, weight filtrations,
and modular level presentations. Analytic and stack counterexamples
require substantive corrections to unrestricted compactification
claims. Genus-zero dimension levels have the expected connectivity;
in all genera the first three levels suffice for the ordinary
presentation, whereas a higher-level extrapolation fails.
\textbf{INCOMPLETE}: the non-toric comparison conjecture for arbitrary
good projective compactifications is not proved or refuted here.
Its morphism case and log-crepancy reduction are proved.
\end{abstract}
\tableofcontents
\clearpage
\section{The ordinary-exit criterion and the two reconstructions}
\label{sec:criterion}
Let $D\subset X$ be a reduced normal-crossings divisor and $U=X\setminus D$.
Write $\mathcal E=\Exit(X,D)$, with arrows pointing from a deeper stratum
to a more open one. At a point $x$ having local branch set $I_x$, put
$\Lambda_x=\Z^{I_x}$ and let $\mu_x:\Lambda_x\to\pi_1(U)$ be the
homomorphism of positive normal meridians, with a chosen connecting path.
Every assertion is componentwise. Neither $\mathcal E$ nor its classifying
space is assumed to be a one-type.

\subsection{Normal links, flags, and specialization}

\begin{theorem}[The normal-link mapping-space formula]\label{P1a1}
\status{Proved}.
Let $S,T$ be distinct incident strata, of depths $d>e$.
There is a finite incidence cover $F_{ST}\to S$ selecting the local
branches that remain zero in $T$, and a torus bundle
\[
 (S^1)^{d-e}\longrightarrow L_{ST}\longrightarrow F_{ST}.
\]
Its endpoint maps give a coherently defined map $L_{ST}\to S\times T$.
For $s\in S,t\in T$,
\[
 \Map_{\mathcal E}(s,t)\simeq
 \hofib_{(s,t)}(L_{ST}\longrightarrow S\times T). \tag{1}
\]
Within one stratum the mapping space is its path space; in particular
$\Map_{\mathcal E}(s,s)\simeq\Omega_sS$.
Nonincident strata have empty crossing mapping spaces.
These assertions hold on stack homotopy types, retaining all inertia.
\end{theorem}
\begin{proof}
In a coordinate chart write $D$ as $z_1\cdots z_d=0$ at $S$.
A local generization to $T$ selects the $e$ coordinates that stay zero.
In the other $d-e$ coordinates a short exit has positive radius and a
nonzero angular coordinate. The radial choices are contractible and
the angular choices form $(S^1)^{d-e}$. Continuation along $S$
permutes branch selections and glues angles by the normal line-bundle
transitions. This produces the incidence cover and torus bundle,
rather than an unjustified globally labelled product.

The holink consists of paths starting in $S$ and immediately entering
$T$. Normal-coordinate contraction identifies its homotopy type with
the angular model, over the starting point and up to homotopy over
the endpoint. Allow arbitrary paths in $T$ to the prescribed endpoint
and paths in $S$ to the prescribed starting point. This is exactly
the homotopy fibre in (1). The holink-to-exit comparison is
\citet[Theorem 4.9]{Miller2009}; normal-crossing charts satisfy its
local tameness hypotheses. For the local-to-global use of its finite
stratification formulation, a compact parameter family of paths is
covered by finitely many such charts and their finitely many local
incidence types. The conical exit construction and
exodromy are those of \citet[Appendix A.9]{LurieHA2017}.

Within a single stratum no direction is exited, so one obtains its
fundamental infinity-groupoid. An exit cannot meet a stratum outside
the incidence closure order. Finally apply these constructions on
every level of an étale atlas and its entire Čech nerve. Branch
covers, links, and flag comparisons commute with étale pullback.
Descent of that coherent diagram proves the stack assertion, as in
\citet[Proposition 3.16 and Corollary 3.17]{JansenToolbox2024}.
No coarse quotient is substituted.
\end{proof}

\begin{theorem}[The normal-bundle and coherent-flag model]\label{P1a2}
\status{Proved}.
On a branch-labelled open incidence stratum $D_J^\circ$, let $N_j$
be its normal lines and put
\[
 P_J^\times=\prod_{j\in J,D_J^\circ}N_j^\times,\qquad
 P_J=P_J^\times/(\R_{>0})^J .
\]
The punctured local tube has homotopy type $P_J^\times\simeq P_J$.
These are the angular interiors of the corresponding closed faces
of $X^{\log}$. Specialization and its coherences are supplied by
closed-face inclusions and all incidence flags.
Let $\mathcal E_D=\Exit(D)$, $\mathcal E_U=\Pi_\infty U$, and
\[
 \mathcal H_X=\Fun(\Delta^1,\mathcal E)
 \times_{\mathcal E\times\mathcal E}
             (\mathcal E_D\times\mathcal E_U). \tag{2}
\]
Its objects are crossing exits and its morphisms are coherent squares.
Then
\[
                  B\mathcal H_X\simeq\partial X^{\log}. \tag{3}
\]
Strict positive-rank flags give a semisimplicial resolution. A full
simplicial resolution also includes repeated ranks and identity
degeneracies. The manuscript's fundamental-groupoid normal diagram
means applying $\Pi_1$ to every full flag space with its coherent
maps; it does not mean forming the crossing category of $h\mathcal E$.
\end{theorem}
\begin{proof}
Choose Hermitian metrics. A nonzero vector has a positive length
and a unit direction, giving
$P_J^\times\simeq P_J\times(\R_{>0})^J$.
The radial factor contracts. A differentiable tubular neighborhood
of the open stratum identifies its punctured normal directions
with the punctured tube. Locally this is the polar-coordinate
description of punctured coordinate discs.

The oriented blow-up replaces a vanishing coordinate by
$(r_j,\theta_j)\in\R_{\ge0}\times S^1$. A chosen branch
intersection defines a closed face. Its angular interior is the
torus bundle above. The interior inclusion is a homotopy
equivalence: compatible collars push the remaining boundary
coordinates slightly positive. Over noncompact strata choose
positive collar functions locally finitely. Forgetting branch
conditions is an actual face inclusion. Transporting it across
these equivalences gives the specialization maps; chains of
inclusions give their homotopies and all higher compatibilities.
Thus inverse equivalences are not chosen as unrelated,
incoherent tubular embeddings.

Alternatively use deformation to the normal cone. The nonzero
deformation fibres are the original tubular pair and the zero
fibre is the normal pair. On punctured normal coordinates the
intervening positive radial parameters are contractible.
Apply this to the entire flag system, then descend through
branch covers and the étale Čech nerve. This gives the
homotopy-level specialization despite the analytic transition
or transvection at scan 146; a global holomorphic embedding
of a normal bundle is not asserted. The logarithmic comparison
is compatible with \citet{KatoNakayama1999}.

In a conical chart a crossing path is determined, up to
contractible radial choices, by its angular direction and
its endpoint paths. Groupoid completion of the endpoint
actions leaves the angular boundary. A chain of squares in
(2) records a chain of incident faces with angular data.
Closed-face inclusions have collars and are cofibrant, so
their homotopy-colimit realization is their union.
The usual Čech nerve is indexed by subsets of faces;
its barycentric flag refinement is the strict-flag
semisimplicial model. Repetitions provide degeneracies.
The chartwise equivalences glue and descend, proving (3).

A loop in a crossing mapping space can be lost in
$h\mathcal E$ although it survives in $\Pi_1$ of the full
link. The examples below exhibit this distinction.
\end{proof}

\subsection{The exact criterion and the failure of adjacent tests}

\begin{theorem}[The ordinary-exit criterion]\label{P1b1}
\status{Proved}. The following are equivalent:
\begin{enumerate}
\item $\mathcal E$ is equivalent to an ordinary category.
\item Every open incidence stratum is aspherical and every full
      $\mu_x:\Lambda_x\to\pi_1(U)$ is injective.
\item Every open incidence stratum is aspherical and every relative
      normal-meridian lattice for every incident pair injects into
      the group of the more open stratum.
\end{enumerate}
Full injection implies all relative injections even without
asphericity. Testing only immediately adjacent strata is insufficient.
\end{theorem}
\begin{proof}
An infinity-category is equivalent to an ordinary category
exactly when its mapping spaces are homotopy-discrete.
Here the endomorphism spaces are $\Omega S$. Their
homotopy-discreteness is equivalent to $\pi_i(S)=0$ for $i\ge2$.

Assume the strata are aspherical. Each finite incidence cover
is aspherical. The torus-bundle homotopy sequence shows that
$L_{ST}$ is aspherical and that its fibre lattice injects.
The kernel of
\[
 \pi_1(L_{ST})\longrightarrow\pi_1(S)\times\pi_1(T)
\]
is exactly the kernel of that fibre lattice in $\pi_1(T)$:
the incidence cover injects on $\pi_1$ into $S$, so an
element in the displayed kernel lies in the fibre lattice.
For a map between aspherical spaces, every component of
its homotopy fibre is a $K(\ker,1)$. Formula (1) is
therefore homotopy-discrete exactly when the relative
kernel vanishes. This proves the third equivalence and
the necessity of the full maps by taking $T=U$.

For the reduction from full to relative maps, write the
source branch set as $I=J\sqcup K$, with $J$ still
vanishing in $T$. Push the full local torus into the
normal torus bundle $P_T\to T$. If $v\in\Z^K$ dies
in $\pi_1(T)$, exactness says its lift in $\pi_1(P_T)$
equals the image of some $w\in\Z^J$. This does not
require $\Z^J\to\pi_1(P_T)$ to inject.
The full vector $v-w$ consequently dies in $\pi_1(U)$.
Full injection implies $v-w=0$, and disjoint supports
give $v=w=0$. No asphericity was used in this reduction.
It is compatible with incidence covers and stack basepoints.
The next example refutes the adjacent shortcut.
\end{proof}

\begin{counterexample}[Adjacent injections with a full kernel]\label{P1b2}
\status{Refuted}.
Even when every stratum is aspherical, injectivity for all adjacent
pairs does not imply goodness.
\end{counterexample}
\begin{proof}
Start with the axes $A=\{x=0\},B=\{y=0\}$ in $\C^2$.
Blow up $p=(0,1)\in A$, write $E$ for the exceptional
curve, and remove a point $r\in E$ different from
$q=E\cap A'$. Use the reduced total boundary.
The open complement remains $U=(\C^*)^2$, and the
boundary curve strata are
\[
 A^\circ=\C\setminus\{0,1\},\quad B^\circ=\C^*,\quad
 E^\circ=\PP^1\setminus\{q,r\}=\C^* .
\]
There are two point strata. All strata are aspherical.
The curve-to-open meridians are $e_x,e_y,e_x$ in $\Z^2$,
so each is injective. At either node, each relative
node-to-curve meridian is a nontrivial peripheral loop
in one of these punctured curves. Its free group is
torsion-free, so the adjacent map injects.
At $q$, however, the full map is
\[
 \Z^2\longrightarrow\Z^2,\qquad (a,b)\longmapsto(a+b)e_x,
\]
with kernel $\Z(1,-1)$. Individual cyclic injection
cannot replace independence of the full lattice.
The integer check is in \texttt{criterion\_checks.py}.
\end{proof}

\begin{proposition}[The missing normal transgression]\label{P1b3}
\status{Proved}.
For a rank-$r$ normal torus bundle over a branch-labelled stratum,
\[
 \pi_2(S)\xrightarrow{\partial}\Z^r
 \longrightarrow\pi_1(P_S)\longrightarrow\pi_1(S)\longrightarrow1.
 \tag{4}
\]
The lattice injects exactly when $\partial=0$.
Its coordinates are the first Chern numbers of the normal
lines on spherical classes. Thus $\pi_2(S)=0$ is sufficient
but not necessary. The lattice quotient may have torsion.
In (37), the base must be the appropriate open incidence
stratum; substituting the closed intersection changes the claim.
\end{proposition}
\begin{proof}
This is the homotopy exact sequence of the torus fibration.
Pull a normal line back to a representative sphere and
trivialize over its hemispheres. Its equatorial clutching
function has degree its first Chern number. Lifting that
equator gives the same integer as the connecting map.
Taking the product gives all coordinates of $\partial$.

For $\mathcal O_{\PP^1}(1)$ the connecting map is
multiplication by one. For a trivial line over $\PP^1$
it is zero although $\pi_2(\PP^1)=\Z$.
A conic in $\PP^2$ has normal degree four, so its
circle-bundle group is $\Z/4$, exhibiting torsion.

For the coordinate axes, the first component's open
stratum is $\C^*$. Its punctured normal tube inside
the open complement has group $\Z^2$, mapping to
$\pi_1(\C^*)=\Z$ with rank-one kernel.
Replacing the base by the closed axis $\C$ gives
zero target and rank-two kernel. The proposed
sequence cannot describe both bases. This supplies
the correction to (37) and its later reminder in
MPL 149, scan 134.
\end{proof}

\subsection{Three operations which must be distinguished}

\begin{theorem}[Derived reconstruction and ordinary truncation]\label{P1c1}
\status{Proved}.
Let $\mathcal T$ be the complete normal-face/link diagram,
refined by open incidence strata and all flags.
\begin{enumerate}
\item Without goodness, its full homotopy types reconstruct $X$
      and its constructible infinity-topos. Its fundamental-groupoid
      diagram always reconstructs $\Pi_1X$, and $h\mathcal E$
      classifies constructible sets.
\item If all open strata are aspherical, the refined values can
      be replaced objectwise by $B\Pi_1$ without changing their
      homotopy colimit. The full ambient homotopy type is recovered
      from this derived groupoid diagram under asphericity alone.
\item With full meridian injection as well,
      $\mathcal E\simeq h\mathcal E$. Both the genuine constructible
      infinity-topos and the infinity-topos associated with
      ordinary constructible sets then have shape $X$.
\end{enumerate}
Item two concerns the refined open-stratum/link diagram;
it does not assert asphericity of closed branch intersections
or identify its homotopy colimit with $B(h\mathcal E)$.
\end{theorem}
\begin{proof}
A normal-crossing collar glues open strata to their normal
directions. Chartwise this is gluing coordinate cones
and faces, which becomes a homotopy pushout using collars.
Iterate through the finite depth and retain all flags.
An open Čech cover and an étale atlas give the global
and stack homotopy-colimit reconstruction.
Equivalently exodromy identifies the constructible
space-valued sheaves with $\Fun(\mathcal E,\mathrm{Spaces})$,
and groupoid completion gives the homotopy type of $X$;
see \citet[Appendix A.9]{LurieHA2017} and
\citet[Proposition 3.24 and Corollary 3.23]{JansenToolbox2024}.
The collar argument supplies the diagram used here.

The fundamental-groupoid functor is left adjoint to
the inclusion of one-types and preserves colimits
formed in groupoids. No asphericity is needed for
the $\Pi_1$ conclusion. Every set-valued functor
on an infinity-category factors uniquely through
its homotopy category. Exodromy therefore also gives
$\Fun(h\mathcal E,\mathrm{Set})$ without goodness.

Under asphericity, refined diagram values are
finite covers of open strata, iterated normal-torus
bundles, and contractible radial parameters.
Their homotopy exact sequences prove asphericity.
The natural maps to $B\Pi_1$ are objectwise
equivalences of the full coherent diagram.
They remain equivalences after homotopy colimits.
Deleting triangular homotopies is a different
operation and can change even $\Pi_1$, as Part 5 shows.

Finally apply Theorem~\ref{P1b1}. For any small category
$\mathcal C$, global sections of a constant space
$K$ in its presheaf infinity-topos are
$\lim_{\mathcal C}K=\Map(B\mathcal C,K)$.
Thus its shape is $B\mathcal C$.
Apply this to $\mathcal E$ and to $h\mathcal E$.
Their shapes agree under goodness. The stronger
ordinary-topos consequence requires this extra
hypothesis and does not follow from item two alone.
\end{proof}

\begin{counterexample}[The separate roles of the two hypotheses]\label{P1c2}
\status{Refuted}.
Asphericity alone does not identify the ordinary constructible
shape with $X$. Full meridian injection alone does not recover
$X$ from its fundamental-groupoid diagram.
\end{counterexample}
\begin{proof}
For $(\PP^1,\{\infty\})$ both strata are contractible
and the link is $S^1$. The derived groupoid diagram
\[
                    *\longleftarrow B\Z\longrightarrow *
\]
has homotopy pushout $S^2$. But the meridian dies
in the open $\C$. The crossing exit space is $S^1$,
so $h\mathcal E=[1]$ and its classifying space is
contractible. The ordinary consequence fails,
although derived reconstruction succeeds.

For the other direction take
\[
             (X,D)=(\PP^1\times\C,\PP^1\times\{0\}).
\]
Its normal meridian maps isomorphically onto
$\pi_1(U)=\Z$, while the closed stratum is
not aspherical. Its closed-link-open groupoid
diagram is
\[
                 *\longleftarrow B\Z\xrightarrow{\id}B\Z ,
\]
with contractible homotopy pushout.
This is exactly the same diagram as for
$(\C^2,\C\times\{0\})$, although the two
ambient homotopy types are $S^2$ and a point.
Thus those groupoid data cannot distinguish them.
\end{proof}

\subsection{Shapes, fibres, and the first possible obstruction}

\begin{theorem}[The two shapes and their gluing data]\label{P1d1}
\status{Proved}.
The genuine constructible infinity-topos has shape
$B\mathcal E\simeq X$. The infinity-topos associated
with the ordinary constructible-set topos has shape
$B(h\mathcal E)$. The canonical comparison is
\[
            q:X\simeq B\mathcal E\longrightarrow B(h\mathcal E). \tag{5}
\]
With two strata, closed $S$ and open $T$, and link $L$,
the genuine shape is
\[
                              S\cup_L^hT. \tag{6}
\]
More strata require every incidence flag.
If $S,T,L$ are aspherical, with groups $A,B,G$
and endpoint homomorphism $G\to A\times B$,
the ordinary crossing homotopy type uses
$\im(G\to A\times B)$ in place of $G$.
\end{theorem}
\begin{proof}
The shape assertions follow from the constant
global-sections calculation above. The closed
stratum has a mapping-cylinder neighborhood with
boundary $L$; gluing to the open complement
gives (6). Shrinking collars does not change
the diagram over its endpoints.

The crossing mapping space is the homotopy
fibre of $BG\to B(A\times B)$. Its components
are cosets of the image, each with group
the kernel. For the ordinary exit category
only the component set remains, acted on
by $A\times B$. The homotopy quotient of
this transitive action is the classifying
space of its stabilizer, the image subgroup.
Before truncation that quotient is $BG$.
This proves the precise change in the
crossing data. More strata require the
same operation coherently over full flags.
\end{proof}

\begin{proposition}[Explicit comparison fibres]\label{P1d2}
\status{Proved}.
The comparison fibres in (5) for
$(\PP^1,\{\infty\})$, $(\PP^n,\PP^{n-1})$,
with $n\ge1$, and the product of two copies of
$(\PP^1,\{\infty\})$ are respectively
\[
                    S^2,\qquad\PP^n,\qquad S^2\times S^2 .
\]
Their first nonzero fibre groups are $\Z,\Z,\Z^2$
in degree two.
\end{proposition}
\begin{proof}
In the first case the ordinary category is $[1]$.
For the hyperplane pair the open is $\C^n$,
the closed stratum has trivial fundamental group,
and the link is the Hopf sphere $S^{2n-1}$.
Its map to $\PP^{n-1}$ is the circle bundle,
so (1) gives crossing mapping space $S^1$;
the closed endomorphism space is
$\Omega\PP^{n-1}$. Passing to components
again gives $[1]$, of contractible shape.
The comparison fibre is therefore $\PP^n$.

Exit categories of product stratifications
are products: paths, exit conditions, and
homotopies are coordinatewise. The third
ordinary category is $[1]\times[1]$,
again contractible, and its genuine shape
is $\PP^1\times\PP^1$.
The cell structures and product formula
give the stated second homotopy groups.
The cellular Betti-number checks are in
\texttt{criterion\_checks.py}.
\end{proof}

\begin{theorem}[Connectivity and the first comparison obstruction]\label{P1d3}
\status{Proved}.
For every small infinity-category $\mathcal E$,
$B\mathcal E\to B(h\mathcal E)$ is two-connected.
If each component of every mapping space is
$(r-1)$-connected, $r\ge1$, it is
$(r+1)$-connected.

Choose a well-based enriched-category model
with discrete objects, and put
\[
 V_p=\coprod_{x_0,\ldots,x_p}
       \Map(x_0,x_1)\times\cdots\times\Map(x_{p-1},x_p),
 \qquad W_p=\pi_0V_p .
\]
Their realizations are $B\mathcal E$ and $B(h\mathcal E)$.
Fix a component $Y$ of the latter. Use its fundamental-groupoid
local system of basepoint paths to form the covering bar
$\widetilde W_\bullet$, and pull this cover back to
$\widetilde V_\bullet$. Their realizations are the universal
cover of $Y$ and its pullback to $B\mathcal E$.
In particular the faces transport lifts along the corresponding
edges; this is not an unrelated levelwise choice of lifts.
For its homotopy fibre $F$,
\[
 \pi_{r+1}F\cong
 \coker\left(
 H_r(\widetilde V_2;\Z)
 \xrightarrow{d_0-d_1+d_2}
 H_r(\widetilde V_1;\Z)\right) . \tag{7}
\]
This is an isomorphism of $\Z[\pi_1Y]$-modules.
It is the first possible obstruction, not necessarily
a nonzero one. The first nonzero disagreement need not
occur in degree two, even for smooth algebraic NC pairs.
\end{theorem}
\begin{proof}
Use a cofibrant simplicial-category model,
so the displayed composition is strict and
the nerve degeneracies are cofibrations.
Each component of $V_p$ is a product of
$(r-1)$-connected spaces, and $V_p\to W_p$
is $r$-connected; at degree zero it is an
isomorphism. In relative realization,
a cell in simplicial degree $p$ acquires
$p$ extra dimensions. The first possible
relative cells in positive simplicial
degree occur in dimension $r+2$, at $p=1$.
The skeletal filtration of the simplicial
mapping cylinder therefore proves
$(r+1)$-connectivity. For $r=1$ this
applies to every mapping-space component.

Over the universal cover of $Y$ the relative
realization spectral sequence is
\[
 E^1_{p,j}=H_j(\widetilde W_p,\widetilde V_p;\Z)
 \Longrightarrow H_{p+j}(\widetilde Y,\widetilde X;\Z).
\]
The degree-zero column vanishes.
For $p>0$ the rows $j\le r$ vanish and
\[
 H_{r+1}(\widetilde W_p,\widetilde V_p)
       \cong H_r(\widetilde V_p).
\]
In total degree $r+2$ the only position
is $(1,r+1)$. Its incoming first
differential is $d_0-d_1+d_2$, and the
outgoing one lands in the zero column.
Higher incoming differentials start
in vanishing lower rows; higher
outgoing ones leave the first quadrant.
The relative homology group is thus
the cokernel in (7).
The already proved connectivity and
relative Hurewicz identify it with
the relative homotopy group in
degree $r+2$, equivalently $\pi_{r+1}F$.
Deck transformations give the module
action throughout. In particular
\[
 \pi_{r+2}(B\mathcal E)\to\pi_{r+2}Y
 \to\pi_{r+1}F\to\pi_{r+1}(B\mathcal E)
 \to\pi_{r+1}Y\to0
\]
is exact. The module measures more
than an unqualified difference of ranks.

For a counterexample to an always-degree-two
claim take $m\ge3$ and
\[
 Q_m=\{\textstyle\sum_{i=0}^m z_i^2=1\}\subset\C^{m+1},
 \qquad (X,D)=(Q_m\times\C,Q_m\times\{0\}).
\]
Writing $z=x+iy$ gives $x\cdot y=0$ and
$|x|^2=1+|y|^2$. The map
$(x,y)\mapsto(x/\sqrt{1+|y|^2},y)$
identifies $Q_m$ with $TS^m$, so
$X\simeq S^m$. The normal meridian
injects isomorphically into
$\pi_1(Q_m\times\C^*)=\Z$.
The ordinary exit category has a
closed object with trivial group,
an open object with group $\Z$,
and a single crossing arrow.
Its closed object is initial.
Thus its classifying space is
contractible, the comparison fibre
is $S^m$, and its first nonzero
group is $\pi_m=\Z$.
This can occur arbitrarily far
above degree two.
\end{proof}


\subsection{Explicit exit categories}
\label{sec:explicit-examples}

All categories below are skeletons: one chosen basepoint represents each
connected stratum. A displayed group is a discrete mapping space unless a
classifying space or loop space is expressly written. Choices of paths and
positive peripheral generators change the descriptions by equivalence.

\begin{proposition}[One puncture]\label{P1e1}
\status{Proved.}
For $(\PP^1,\{\infty\})$ the exit category has objects $p,u$,
\[
 \Map(p,p)=\Map(u,u)=*,\qquad \Map(p,u)\simeq S^1,
 \qquad \Map(u,p)=\varnothing.
\]
The only compositions are unit compositions. Its homotopy category is
$[1]$, whereas its classifying space is $S^2$.
\end{proposition}
\begin{proof}
Both strata are contractible. The normal link of $p$ is a circle, and both
endpoint maps have contractible targets. Proposition~\ref{P1a1} therefore
identifies the crossing mapping space with this circle. There are no two
composable nonidentity crossing arrows. The realization of this enriched
two-object category is the unreduced suspension of $S^1$, hence $S^2$;
on taking path components of the mapping spaces it becomes the interval.
Equivalently, the normal meridian maps $\Z\to\pi_1(\C)=1$, with kernel
$\Z$. This is the smallest example where condition (i) holds and (ii)
fails.
\end{proof}

\begin{proposition}[Three punctures]\label{P1e2}
\status{Proved.}
For $(\PP^1,\{0,1,\infty\})$, put
\[
 G=F(a,b),\qquad H_0=\langle a\rangle,\quad
 H_1=\langle b\rangle,\quad H_\infty=\langle(ab)^{-1}\rangle.
\]
The exit category is ordinary. Its four objects are $u,p_0,p_1,p_\infty$;
$\Aut(u)=G$, $\Aut(p_i)=1$, and
\[
 \Hom(p_i,u)=G/H_i.
\]
There are no other off-diagonal arrows. Postcomposition at $u$ is left
multiplication by $G$. The classifying space of this ordinary category is
$\PP^1\simeq S^2$.
\end{proposition}
\begin{proof}
The punctured sphere deformation retracts to a wedge of two circles.
Positive small loops about the punctures can be based so that their
product is $1$, giving the three displayed peripheral subgroups. Each
peripheral word is nontrivial and therefore has infinite order in the
free group. The normal-lattice maps are injective and all strata are
aspherical, so Proposition~\ref{P1b1} applies. The homotopy fibre of
$B\Z\to BG$ consists of contractible components indexed by the left
$G$-set of left cosets $G/H_i$ (the quotient by the right $H_i$-action).
This proves the arrow and composition
formulas. Derived reconstruction then identifies the realization with
the sphere. In particular, goodness does not imply that the ambient
compactification is aspherical.
\end{proof}

\begin{proposition}[The stable five-pointed rational curve]\label{P1e3}
\status{Proved.}
The ordinary exit category of
$(\overline{\mathcal M}_{0,5},\partial\overline{\mathcal M}_{0,5})$
has $26$ objects: one open object, ten divisor objects $d_I$ indexed by
two-element subsets $I\subset\{1,\ldots,5\}$, and fifteen node objects
$q_{I,J}$ indexed by unordered disjoint pairs of such subsets. Its
automorphism groups are, respectively,
\[
 G_5=\Gamma_{0,5},\qquad F_2,\qquad 1.
\]
Choose a multicurve $\sigma_s$ representing each stratum, with
$\sigma_u=\varnothing$, and let $\Delta_\sigma$ be the subgroup generated
by the positive Dehn twists about its components. A complete formula is
\begin{equation}\label{eq:M05-transporter}
 \Hom(s,t)=
 \{g\in\Gamma_{0,5}:g\sigma_s\supseteq\sigma_t\}/\Delta_{\sigma_s},
 \qquad [h]\circ[g]=[hg].
\end{equation}
Here the quotient is by right multiplication. In particular,
\[
 \Hom(d_I,u)=G_5/\langle T_I\rangle,\qquad
 \Hom(q_{I,J},u)=G_5/\langle T_{I,q},T_{J,q}\rangle,
\]
where $T_{I,q},T_{J,q}$ are the two compatible twists in the chosen
node multicurve $\sigma_q$. They need not equal the independently
chosen divisor representatives $T_I,T_J$; the transporter formula
includes the required conjugations.
The arrows from a node to either incident divisor form
$F_2/\langle c\rangle$, where $c$ is that puncture's peripheral loop.
There are $30$ node--divisor incidences. All other off-diagonal arrow
sets are empty.
\end{proposition}
\begin{proof}
Fixing three marked points at $0,1,\infty$ identifies the open moduli
space with $\operatorname{Conf}_2(\C\setminus\{0,1\})$.
Forgetting the second moving point is a locally trivial fibration with
base $\C\setminus\{0,1\}$ and fibre a three-punctured affine line.
Both are aspherical, and the homotopy exact sequence proves that the
total space is aspherical. The divisor strata are copies of
$\mathcal M_{0,4}$, and the zero-dimensional strata are points.
The local plumbing equation $xy=t$ identifies a positive normal
meridian with the positive Dehn twist about the vanishing curve.
Twists about disjoint essential curves generate a free abelian group:
for each curve choose a test curve crossing it and disjoint from the
other multicurve components; cutting along an annular neighbourhood
shows that a nonzero power changes its winding number. A product of
twists acting trivially would therefore have every exponent zero.
Thus all full normal lattices inject, and the criterion proves that
the exit category is ordinary.

The augmented Teichm\"uller description now gives
\eqref{eq:M05-transporter}: a marking identifies the source multicurve
with a multicurve containing the target one, and changing the marking
by twists around collapsed source curves has no effect. This is the
transport-category construction of \citet[Construction~3.18,
Corollary~6.6 and Theorem~6.7]{Jansen2023}. It can also be read directly
from the normal torus charts and their marking changes. The formula
specifies the full category, not only its cardinalities. It is
well-defined: if $h$ is replaced by $h\delta_t$, then
$g^{-1}\delta_tg\in\Delta_{\sigma_s}$ because
$\sigma_t\subseteq g\sigma_s$. Replacement of $g$ by a source twist is
already absorbed by the right quotient. Composition is therefore
associative and the identity representatives are units.

The diagonal case is
$\Stab(\sigma_s)/\Delta_{\sigma_s}$, which is $G_5$, $F_2$, or $1$
according to the three stratum types. A divisor is specified by a stable
partition $2+3$ of the five labels, giving $\binom52=10$ choices. A
maximal two-edge tree is specified by two disjoint two-element subsets,
giving $\frac12\binom52\binom32=15$ choices. Each such node lies on
exactly two divisors. The presentation of $G_5$ and the complete finite
verification of these incidence counts are supplied below and in the
code archive.
\end{proof}

\begin{proposition}[Smooth toric surfaces]\label{P1e4}
\status{Proved.}
Let $X_\Sigma$ be a smooth complex toric surface with fan $\Sigma$ in
$N\cong\Z^2$, and let $D$ be its toric boundary. For a cone $\sigma$,
put $N_\sigma=N\cap\operatorname{span}_{\R}(\sigma)$.
The ordinary exit category has one object for each cone and
\begin{equation}\label{eq:toric-exit}
 \Hom(\sigma,\tau)=
 \begin{cases}
 N/N_\sigma,&\tau\text{ is a face of }\sigma,\\
 \varnothing,&\text{otherwise}.
 \end{cases}
\end{equation}
In particular $\Aut(\sigma)=N/N_\sigma$. For
$\upsilon\subseteq\tau\subseteq\sigma$, composition sends
\[
 (a,b)\longmapsto a+\operatorname{pr}_{\tau,\sigma}(b),
 \quad a\in N/N_\sigma,\quad b\in N/N_\tau.
\]
The pairs $(\PP^2,D_{\rm tor})$ and
$(\mathbb F_n,D_{\rm tor})$ have respectively seven and nine objects.
\end{proposition}
\begin{proof}
The orbit $O_\sigma$ is $(\C^*)^{2-\dim\sigma}$, with fundamental group
$N/N_\sigma$. A smooth cone is generated by part of an integral basis,
so its normal-meridian map is the inclusion of its free lattice into
$N$. Thus the criterion applies to every cone. If $\tau$ is a face of
$\sigma$, the angular incidence space for $O_\sigma$ to $O_\tau$ has
fundamental group $N/N_\tau$. Its endpoint homomorphism is
\[
 N/N_\tau\longrightarrow (N/N_\sigma)\times(N/N_\tau),
 \qquad v\longmapsto(\operatorname{pr}_{\tau,\sigma}v,v).
\]
It is injective, and the components of its homotopy fibre are naturally
$N/N_\sigma$. This proves \eqref{eq:toric-exit}. Compatible choices of
torus basepoints identify concatenation with addition and projection,
giving the stated composition. When $\tau$ is not a face, the closure
of $O_\tau$ does not contain $O_\sigma$, so no exit is possible.
The complete fan of $\PP^2$ has one zero cone, three rays and three
two-cones; that of $\mathbb F_n$ has one, four and four.
\end{proof}

\begin{proposition}[Products of marked curves]\label{P1e5}
\status{Proved.}
Let $X=\prod_{i=1}^m C_i$ be a product of connected smooth compact
curves, with $g_i=\operatorname{genus}(C_i)$, and let
\[
 D=\bigcup_i\operatorname{pr}_i^{-1}(D_i),\qquad r_i=|D_i|.
\]
The full exit infinity-category is the product of the curve exit
categories. For $r_i>0$, the $i$th factor has the open object $u_i$ and
one object for each marked point; its open automorphism group is
$G_i=F_{2g_i+r_i-1}$, and
\begin{equation}\label{eq:curve-exit-general}
 \Map(p,u_i)\simeq
 \coprod_{G_i/\im\mu_p} B\ker(\mu_p:\Z\to G_i).
\end{equation}
The point endomorphism spaces are contractible, and all other
off-diagonal mapping spaces are empty. If $r_i=0$, the factor is
$\Pi_\infty C_i$; its one-object model has endomorphism space
$\Omega C_i$.

Condition (i) holds exactly when no factor has $(g_i,r_i)=(0,0)$.
Condition (ii) holds exactly when no factor has $(g_i,r_i)=(0,1)$.
Consequently the product exit category is ordinary exactly when both
types are absent. Formula~\eqref{eq:curve-exit-general}, products of
mapping spaces, and coordinatewise compositions give all remaining
cases, including the nonordinary ones.
\end{proposition}
\begin{proof}
A punctured compact curve retracts to a graph of rank $2g+r-1$.
The peripheral word at a puncture has infinite order except for the
one-punctured sphere. For at least two punctures this follows from the
standard free presentation; with one puncture and $g\ge1$ it is the
nontrivial reduced product of the $g$ handle commutators. Free groups
are torsion-free. The normal link is $B\Z$, so its homotopy fibre over
the aspherical open curve is precisely
\eqref{eq:curve-exit-general}. The exceptional one-punctured sphere
therefore has crossing mapping space $S^1$. Postcomposition is the
left $G_i$-action on the coset components, with its induced action on
the fibre; this follows from path lifting and fixes the composition
data in the nonordinary case as well.

Exit paths and their higher homotopies in a product are specified
coordinatewise. Allowing a coordinate to pause lets any finite set of
coordinate exit times share a parametrization, so this identification
gives the product infinity-category. The strata are products of open
curves and marked points. An unmarked sphere factor has nonzero
$\pi_2$ in every stratum; every other curve factor is aspherical on
each of its strata. This proves the assertion about (i). Each full
normal lattice has one generator for each coordinate constrained to
a marked point, and the images lie in different factors of
$\pi_1(\prod_i(C_i\setminus D_i))$. Such a product homomorphism is
injective exactly when each occurring peripheral homomorphism is
injective, proving the assertion about (ii). An unmarked sphere has
no normal coordinate and hence does not itself violate (ii).
\end{proof}

\subsection{The cuspidal fibre and finite monodromy}
\label{sec:cusp}

\begin{proposition}[Rooting a fixed boundary cannot repair a lattice kernel;
\status{Proved}]\label{P1e-roots}
Suppose the normal-meridian map $\rho:\Z^r\to\pi _1(U)$ of a
normal-crossing pair has nonzero kernel. Taking arbitrary finite roots
of those same boundary components does not make this map injective
on a root-stack atlas. The assertion holds both discretely and
profinally.
\end{proposition}
\begin{proof}
For root indices $n_1,\ldots,n_r$, an atlas has coarse coordinates
$x_i=z_i^{n_i}$. Its meridian map is
$\rho\circ\operatorname{diag}(n_1,\ldots,n_r)$. The subgroup
$\bigoplus_i n_i\Z$ has finite index in $\Z^r$ and meets every
nonzero subgroup nontrivially: multiply any nonzero kernel vector
by a common multiple of the $n_i$. Thus the new map still has a
nonzero integral kernel. The same integral vector is nonzero in
$\widehat\Z^r$, and gives the profinite assertion. More generally,
if a purely profinite kernel vector $v$ is nonzero, multiplication
by a common root index $N$ puts it in the root lattice and cannot
kill it: $\widehat\Z^r$ has no nonzero integral torsion. This is a statement
about rooting an unchanged resolved boundary, not about replacing
the degeneration by a different stable model.
\end{proof}

We distinguish the affine surface $y^2=x^3+t$ from its projective
Weierstrass completion in the fibre direction. Both interpretations
of the cuspidal-fibre example can be computed completely.

\begin{proposition}[The resolved affine cusp; \status{Proved}]
\label{P1e-cusp}
Let $f(x,y)=y^2-x^3$, let $U=\C^2\setminus f^{-1}(0)$, and resolve
the reduced cusp divisor minimally to normal crossings. The resulting
pair has eight connected strata: the open $U$, four curve strata
$C^\circ,E_1^\circ,E_2^\circ,E_3^\circ$, and three nodes
$\nu_i=C_i\cap E_3$, where $C_1=E_1,C_2=E_2,C_C=C$.
They are
\[
 C^\circ=\C^*,\qquad E_1^\circ=E_2^\circ=\mathbb A^1,\qquad
 E_3^\circ=\PP^1\setminus\{0,1,\infty\}.
\]
Every stratum is aspherical. Put
\[
 G=\pi _1(U)=\langle a,b\mid a^2=b^3\rangle,\qquad
 h=a^2=b^3,\qquad m=ab^{-1}=h(ba)^{-1}.
\]
The meridians of $C,E_1,E_2,E_3$ are respectively $m,b,a,h$.
All four have infinite order. The full node maps have kernels
\[
 \ker\bigl(\Z\langle b\rangle\oplus\Z\langle h\rangle\to G\bigr)
       =\Z(-3,1),\qquad
 \ker\bigl(\Z\langle a\rangle\oplus\Z\langle h\rangle\to G\bigr)
       =\Z(-2,1),
\]
whereas $\Z\langle m\rangle\oplus\Z\langle h\rangle\to G$ is injective.
Thus the ordinary-exit-category criterion fails at exactly the two
exceptional-tail nodes.
\end{proposition}
\begin{proof}
The three relevant blow-up charts give, successively,
\[
 f=x^2(v^2-x),\qquad
 f=u^2v^3(v-u),\qquad
 f=u^6w^3(w-1).
\]
The last exceptional component meets the strict cusp and the two
earlier exceptionals at three distinct points. The multiplicities
of the total fibre on $(C,E_1,E_2,E_3)$ are $(1,2,3,6)$, and the
exceptional self-intersections are $(-3,-2,-1)$. The strict affine
cusp is its normalization with the origin removed, hence $\C^*$;
each early exceptional has just one puncture.

The weighted radial action $(x,y)\mapsto(r^2x,r^3y)$ identifies the
complement's homotopy type with the complement of the $(2,3)$ torus
knot. Splitting the knot exterior into its two solid-torus pieces
and applying van Kampen identifies their core powers $a^2$ and
$b^3$, giving the displayed presentation. Equivalently, the base
change $t=s^6$, $x=s^2X$, $y=s^3Y$ trivializes the fibration as
\[
 F\times\C^*,\qquad F=\{Y^2=X^3+1\}.
\]
The smooth projective completion of $F$ is elliptic and has one
point at infinity. Its affine fundamental group is free of rank
two, and the fibration over $\C^*$ proves that $U$ is aspherical.
The degree map induced by $f$ is
\[
 \varphi:G\to\Z,\qquad
 \varphi(a)=3,\quad\varphi(b)=2,\quad\varphi(h)=6,\quad\varphi(m)=1.
\]
The meridian identification can also be read without any ambiguity
from the resolution plumbing. A rational tail of square $-d$ imposes
$\mu_{\rm adjacent}=\mu_{\rm tail}^{d}$: this is the Euler-number
relation in the circle bundle of its normal line. Thus the two tails
impose $b^3=h$ and $a^2=h$. The central component has square $-1$,
so its three peripheral meridians satisfy $bam=h$, with a compatible
cyclic order and normal framing. Hence $m=h(ba)^{-1}=ab^{-1}$.
This identifies the chosen presentation and all four meridians
simultaneously. Their displayed degrees prove that all are infinite.
At the two tail nodes the images are the cyclic groups generated by
$b$ and $a$, with the exact kernels stated above. At the cusp node,
pass to $G/\langle h\rangle=C_2*C_3$. The image of $m=ab^{-1}$ is a
cyclically reduced word of positive alternating length, so no
nonzero power is trivial. A relation $m^rh^s=1$ consequently has
$r=0$, and then $s=0$ by $\varphi(h)=6$. The curve and point
strata have the claimed asphericity directly.
\end{proof}

\paragraph{All objects, mapping spaces, actions, and compositions.}
Here is a concrete group-and-link model for the entire exit category,
not merely for its nontruncated homotopy groups. There is one object
for each of the eight connected strata. The endomorphism groups
in the homotopy category are
\[
 G,\quad \Z\ (C^\circ),\quad 1\ (E_1^\circ),\quad
 1\ (E_2^\circ),\quad F_2\ (E_3^\circ),\quad 1,1,1.
\]
These are the components of the endomorphism mapping spaces; those
spaces themselves are homotopy-discrete. For each curve stratum
$S$ with meridian $\mu_S$, its exit space to $U$ is the discrete
coset set $G/\langle\mu_S\rangle$. The open group acts on the left.
For $C^\circ$ choose the normal framing in which the generator
of its stratum group lifts to $h$; this specifies its right action
on $G/\langle m\rangle$. For $E_3^\circ$ write
\[
 F_2=\langle p_1,p_2\rangle,\qquad p_3=(p_1p_2)^{-1}.
\]
The right action on $G/\langle h\rangle$ is induced by
$p_1\mapsto b,\ p_2\mapsto a$. The tail source groups are trivial.

The remaining nonempty off-diagonal spaces are
\[
\begin{array}{c|c}
\text{exit space}&\text{homotopy type}\\ \hline
\Map(\nu_1,U)&\displaystyle\coprod_{G/\langle b\rangle}S^1\\
\Map(\nu_2,U)&\displaystyle\coprod_{G/\langle a\rangle}S^1\\
\Map(\nu_C,U)&G/\langle m,h\rangle\\
\Map(\nu_1,E_1^\circ),\ \Map(\nu_2,E_2^\circ)&S^1\\
\Map(\nu_C,C^\circ)&*\\
\Map(\nu_i,E_3^\circ)&F_2/\langle p_i\rangle\quad(i=1,2,3).
\end{array}
\]
All other off-diagonal spaces are empty. To specify every face map
and composition, use the following link-group diagram and take
nerves and homotopy fibres, as in the link reconstruction in Part~1.
The punctured normal-link groups of the curves and their maps to $G$
are
\[
\begin{array}{c|c|c}
S&\pi _1(P_S)&\pi _1(P_S)\to G\\ \hline
E_1^\circ&\Z\langle b\rangle&b\mapsto b\\
E_2^\circ&\Z\langle a\rangle&a\mapsto a\\
C^\circ&\Z\langle m,h\rangle&(m,h)\mapsto(m,h)\\
E_3^\circ&\Z\langle h\rangle\times F_2&
             (h,p_1,p_2)\mapsto(h,b,a).
\end{array}
\]
At a node take its full rank-two group, with ordered generators
$(\mu_{\rm leaf},h)$. Its map to the leaf link is respectively
\[
 (b,h)\mapsto(b,b^3),\qquad
 (a,h)\mapsto(a,a^2),\qquad
 (m,h)\mapsto(m,h).
\]
Its map to the central link sends these generators respectively to
\[
 (p_1,h),\qquad(p_2,h),\qquad(hp_3,h).
\]
Projection from a curve link to its stratum kills its normal
generator. These formulas commute with the displayed maps to $G$;
in particular $hp_3\mapsto h(ba)^{-1}=m$. The nerves, their
homotopy fibres, and the induced maps for the two routes from a
node to $U$ specify a simplicial model of the category, including
composition homotopies. The tail-circle generator maps to
$(-3,1)$ or $(-2,1)$ in the corresponding node-to-open space.
Thus neither these circles nor their composition data may be
discarded. The displayed mapping spaces also follow directly from
the homotopy-fibre formula: for a homomorphism of discrete groups,
each component of its classifying-space fibre is the classifying
space of its kernel.

\begin{proposition}[The projective Weierstrass interpretation;
\status{Proved}]\label{P1e-cusp-proper}
For the projective family $y^2z=x^3+tz^3$, the same resolution has
$C^\circ=\mathbb A^1$ instead of $\C^*$; all three leaves are affine lines.
Its open fundamental group is
\[
 G_{\rm pr}=\Z^2\rtimes_A\Z,\qquad
 A=\begin{pmatrix}0&-1\\1&1\end{pmatrix},\quad A^6=1.
\]
It is the quotient of $G$ by the normal closure of
$\lambda=hm^{-6}$. The three node kernels, in the leaf-central
ordering, are now $\Z(-3,1)$, $\Z(-2,1)$, and $\Z(-6,1)$.
Every stratum is still aspherical, but all three full node lattices
fail injectivity. In the previous category replace $G$ by
$G_{\rm pr}$, replace the $C^\circ$ endomorphism group by $1$,
replace $\Map(\nu_C,C^\circ)=*$ by $S^1$, and replace
$\Map(\nu_C,U)$ by $\coprod_{G_{\rm pr}/\langle m\rangle}S^1$.
The third leaf-link homomorphism becomes
$(m,h)\mapsto(m,m^6)$. These replacements, with the same central
link diagram, give the entire projective exit category.
\end{proposition}
\begin{proof}
Filling the unique puncture of $F$ kills its peripheral commutator,
which is the knot longitude $\lambda=hm^{-6}$. Indeed the framing
generator $h$ goes once along the knot and has linking number six,
while $m$ is the positive meridian of linking number one; subtracting
six meridians gives the primitive zero-linking longitude. The remaining
fibre is a torus. The section at infinity splits the fundamental
group sequence, and its order-six monodromy acts on a suitable
homology basis by $A$. For an explicit compatible realization in
the semidirect product, with product
$(v,n)(w,k)=(v+A^nw,n+k)$, take
\[
 a=(e_2,3),\qquad b=(e_1,2),\qquad
 m=(0,1),\qquad h=(0,6).
\]
They satisfy $a^2=b^3=m^6=h$ and generate the group, since $a,b$
together with $m$ generate the two fibre basis vectors. The degree
map is still the second coordinate, so the three leaf meridians
are infinite, and the stated kernels follow from the three exact
power relations. The local resolution is unchanged. The proper
strict transform has square $-6$; the full fibre intersection
matrix annihilates the multiplicity vector $(1,2,3,6)$.
The torus fibration over the punctured disc is aspherical, and
the remaining strata were already aspherical. The homotopy-fibre
calculation of the preceding paragraph gives all the replacements.
\end{proof}

\begin{proposition}[The good orbifold replacement; \status{Proved}]
\label{P1e-cusp-stack}
Let $\zeta$ be a primitive sixth root of unity and let
\[
 \mathcal X=[F\times\Delta_s/\mu _6],\qquad
 \zeta\cdot(X,Y,s)=(\zeta^{-2}X,\zeta^{-3}Y,\zeta s).
\]
Its boundary is $\mathcal F=[F/\mu _6]$ at $s=0$, and its open
complement is the original affine family over $\Delta_t^*$,
where $t=s^6$, $x=s^2X$, and $y=s^3Y$.
The pair is good. Its boundary is the orbifold affine line with
cone orders $2,3$, so
\[
 \pi _1(\mathcal F)=C_2*C_3=G/\langle h\rangle .
\]
Its normal meridian is $h$. The exit category is ordinary up to
equivalence, with two objects, automorphism groups
$C_2*C_3$ and $G$, and crossing set $G/\langle h\rangle$.
The right action is that of the boundary quotient and the left
action is that of $G$.

Replacing $F$ by its projective elliptic completion gives the
corresponding good projective orbifold model. Its boundary group
is $\Z^2\rtimes_A\mu _6=G_{\rm pr}/\langle h\rangle$; its coarse
boundary is $\PP^1$ with cone orders $2,3,6$. The same two-object
description holds with these groups. This replacement is not
obtained by merely rooting the fixed three-blow-up NC divisor.
\end{proposition}
\begin{proof}
The action is free for $s\ne0$. The displayed invariant formulas
identify its quotient there with the given family. The affine
quotient coordinate on the boundary is $w=Y^2=X^3+1$.
At $w=0$ the stabilizer has order two; at $w=1$ it has order
three. The projective completion adds the fixed point at infinity,
of order six. The finite atlases $F$ and $\overline F$ are
aspherical, so their orbifold classifying spaces are aspherical
as well: the Borel fibrations over $B\mu _6$ have no higher
homotopy groups. A circle in the atlas normal coordinate $s$
is the central fibre of this punctured normal circle bundle.
In the affine group the centre is $\langle h\rangle$, since
its quotient $C_2*C_3$ is centreless. In the proper group the
centre is $\langle(0,6)\rangle$, since $A$ has no nonzero fixed
vector. The circle has degree six in the $t$-base, so it is
exactly $h$ and has infinite order. Its full profinite completion
also injects, since the composite with the completed degree map is
multiplication by six on $\widehat\Z$.

The punctured normal-link group is $G$ (respectively $G_{\rm pr}$);
its map to the boundary group is quotient by $\langle h\rangle$,
and its map to the open group is the identity. Taking the
homotopy fibre of the resulting graph homomorphism gives the
discrete crossing set asserted above, with precisely the stated
actions and group multiplication compositions. The criterion
from Part~1 applies.

The affine family is of stable pointed type $(1,1)$ after adding
its missing point. In the projective interpretation that point
can remain a marking rather than a removed boundary component.
We do not call unpointed genus one a stable Deligne--Mumford
curve. The construction is a base change, replacement by the
smooth elliptic model, and quotient descent. Proposition
\ref{P1e-roots} proves separately why finite roots of the old
resolved boundary cannot have the same effect.
\end{proof}

\begin{proposition}[What finite monodromy does and does not imply;
\status{Proved}]\label{P1e-finite}
Finite nontrivial monodromy does not by itself force failure of
meridian injectivity. For example, let a finite cyclic group act
freely on an aspherical smooth curve $F$ and act on a disc by
rotation. Then $(F\times\Delta)/\mu_n$ is a smooth manifold,
its boundary is the aspherical curve $F/\mu_n$, and the normal
meridian has degree $n$ in the punctured-disc base, so the pair
is good. The cusp obstruction above is the exact dependence
among several resolved meridians. Stable orbifold replacement
retains the finite monodromy in the boundary stabilizers while
retaining the infinite normal meridian in the open fundamental
group. All assertions here are over $\C$ or in characteristic
zero; no claim about wild reduction at $2$ or $3$ is made.
\end{proposition}
\begin{proof}
Freeness on $F$ makes the diagonal quotient a manifold even at
$s=0$. The finite covers $F\to F/\mu_n$ and
$F\times\Delta^*\to (F\times\Delta^*)/\mu_n$
identify their higher homotopy groups with those of aspherical
spaces, so both strata are aspherical. In a local normal coordinate
the invariant base parameter is $t=s^n$; the normal circle has degree
$n$ under the map of the open complement to $\Delta_t^*$.
Its cyclic fundamental group therefore injects. The remaining
statements follow from the preceding three cusp calculations and
the fixed-boundary root obstruction.
\end{proof}

\section{Blow-ups and the normal-meridian criterion}
\label{sec:blowups}

Call a normal-crossing pair $(X,D)$ \emph{good} when every connected
stratum, including $U=X\setminus D$, is aspherical and every full
normal-meridian lattice $\Lambda_S\to\pi _1(U)$ is injective. Branches
are taken on the normalization; a self-intersection contributes two
local branches. A \emph{stratal centre} is a smooth closed centre
which, in every normal-crossing chart, is the intersection of a
subset of the boundary branches. Disjoint unions of such centres
are allowed. This local condition is more precise than demanding
that the centre be one globally irreducible stratum closure.

\begin{theorem}[The exact blow-up criterion; \status{Proved}]
\label{P2a}
Let $(X,D)$ be good, and let $Z\subset D$ be a smooth centre having
normal crossings with $D$. Write $\beta:\widetilde X=\operatorname{Bl}_Z X
\to X$ and $\widetilde D=(\beta^{-1}D)_{\rm red}$. Then
$(\widetilde X,\widetilde D)$ is good if and only if $Z$ is stratal.
A nonstratal blow-up creates an explicit integral kernel in a full
normal-meridian lattice, even when the individual new meridian is
infinite. The same kernel obstructs profinite injectivity.
\end{theorem}
\begin{proof}
At a general point of a connected component of the centre choose
coordinates in which
\[
 D=\{x_1\cdots x_s=0\},\qquad
 Z=\{x_1=\cdots=x_s=y_1=\cdots=y_t=0\}.
\]
The normal-crossing hypothesis on the centre permits precisely this
form. The centre is stratal exactly when $t=0$. Denote the old
meridians by $m_1,\ldots,m_s$. The exceptional meridian is
$m_E=m_1\cdots m_s$: a transverse exceptional circle winds once
around every $x_i$, and zero times around coordinates not defining
the boundary.

If $t>0$, use the blow-up chart with exceptional coordinate $y_1$,
where $x_i=y_1u_i$. The exceptional component and all $s$ strict
boundary transforms meet along $y_1=u_1=\cdots=u_s=0$. Their
rank-$(s+1)$ normal lattice maps into the old rank-$s$ lattice by
\[
 (a_1,\ldots,a_s,b)\longmapsto(a_1+b,\ldots,a_s+b).
\]
Its kernel is the primitive vector $(-1,\ldots,-1,1)$, so the
composite map to $\pi _1(U)$ cannot be injective. Its image in
$\widehat\Z^{s+1}$ is still nonzero. This argument does not require
the old pair to be good, and does not assert that $m_E$ itself dies.

If $t=0$, in the chart indexed by $j$ the exceptional ray replaces
$e_j$ by $e_1+\cdots+e_s$, and the other old coordinate rays remain.
This is a unimodular change of basis. Every smaller new cone gives
a primitive sublattice of that same old lattice. Consequently all
new full normal lattices inject into $\pi _1(U)$. Over an old
stratum, each new stratum is, after a finite cover labelling the
branches, a bundle with fibre a complex torus $(\C^*)^a$. The
fibration homotopy sequence proves asphericity from that of the old
stratum. Finite covers do not alter higher homotopy groups.
Strata off the centre are unchanged. Thus the new pair is good.
\end{proof}

\begin{theorem}[Functoriality, ends, and invariance in both directions;
\status{Proved}]\label{P2b}
The blow-down induces a functor
$\Exit(\widetilde X,\widetilde D)\to\Exit(X,D)$ and is the identity
on the open stratum $U$. It preserves the pro-homotopy type of the
punctured neighbourhood of the indicated boundary. For compact
boundary and ordinary collar links, it preserves the homotopy type
of that link. For a stratal blow-up, goodness holds on one side if
and only if it holds on the other. These assertions concern the
boundary end, not unspecified additional ends of a nonproper $X$.
\end{theorem}
\begin{proof}
Each new stratum maps into an old stratum, with incidence order
preserved. Postcomposing exit paths and their simplices with
$\beta$ therefore gives the stated functor. The inverse image of
$U$ is identified with $U$, so its restriction there is the identity.

For the assertion about ends, let $V$ run through neighbourhoods
of $D$ in $X$. The punctured inverse images are literally
\[
 \beta^{-1}(V)\setminus\widetilde D=V\setminus D.
\]
They are cofinal among neighbourhoods of $\widetilde D$. Indeed,
if $W$ contains $\widetilde D$, then the proper map $\beta$ sends
the closed set $\widetilde X\setminus W$ to a closed set disjoint
from $D$; its complement is a suitable $V$. The punctured
neighbourhood pro-spaces thus agree. When a compact boundary has
a collar link, this system is represented up to homotopy by its
link, giving the additional statement.

Forward preservation was proved in Theorem \ref{P2a}. For the
reverse implication suppose the stratal blow-up is good. At a
point of an old stratum with full lattice $\Z^r$, a maximal new
cone has a unimodular rank-$r$ chart. Its injectivity consequently
proves injectivity of the old full lattice.

It remains to recover asphericity of an old stratum $S$ contained
in the centre. Pass to its finite incidence cover, if needed, to
label all $r$ normal branches and the $s$ branches defining the
centre. Let $P_S\to S$ be the punctured full normal bundle, with
fibre $(\C^*)^r$. Its homotopy sequence has connecting map
\[
 \partial:\pi _2(S)\longrightarrow\Z^r.
\]
Every element in its image dies in $\pi _1(P_S)$ and hence in
$\pi _1(U)$. The just-proved full-lattice injection forces
$\partial=0$. The open exceptional stratum $Q$ over $S$ is the
torus bundle with fibre
$(\C^*)^s/\C^*_{\rm diagonal}\cong(\C^*)^{s-1}$.
Its connecting map is the projection of $\partial$ to
$\Z^s/\Z(1,\ldots,1)$, and is therefore zero. Since $Q$ is a
stratum of the good blown-up pair, it is aspherical. The homotopy
sequence of $Q\to S$ now gives $\pi _2(S)=0$, and gives
$\pi_j(S)=0$ for $j\geq3$ because the fibre is aspherical.
This proves the assertion after descent from the incidence cover.
In particular, this argument never uses a nonexistent natural
map $Q\to U$; it uses the punctured normal bundle $P_S$.
\end{proof}

\section{Good compactifications}
\label{sec:compactifications}

\subsection{Surfaces: an exact criterion in a fixed meromorphic class}

A marked completion is a smooth compact complex surface $X$ with
reduced normal-crossing boundary $D$ and a specified isomorphism
$X\setminus D\cong U$. In the algebraic version all these data are
algebraic. Isomorphisms of marked completions induce the identity
of $U$. Two analytic marked completions belong to the same
\emph{meromorphic class} when that identity extends bimeromorphically.
This qualification is essential: a biholomorphism of noncompact
analytic surfaces need not extend meromorphically to completions.

The weighted incidence graph has one vertex for every normalized
irreducible boundary curve, decorated by its genus and
self-intersection, and one edge for every node. A self-node is a
loop, counted twice in the valency. Its \emph{core} is obtained
by repeatedly deleting genus-zero vertices of valency at most
one. Positive-genus vertices are never deleted. The complement
of the core is the \emph{fringe}; when the core is nonempty in
a connected component, its fringe consists of rational trees,
each attached to that core at one point. A \emph{corner blow-up}
is the blow-up of a boundary node, including a self-node.

\begin{lemma}[Good boundaries have no rational fringe;
\status{Proved}]\label{P3a-core-necessity}
The boundary graph of a good surface pair equals its core.
\end{lemma}
\begin{proof}
An isolated smooth rational component is a $\PP^1$ stratum and
is not aspherical. If a rational component has just one incident
branch, its open stratum is $\mathbb A^1$. Let $m$ be its normal
meridian and let $n$ be the adjacent branch meridian at that node.
A punctured normal bundle over $\mathbb A^1$ has fundamental
group $\Z\langle m\rangle$. Hence the two-dimensional node
lattice maps into that cyclic group before mapping to
$\pi _1(U)$, and has a nonzero kernel. Thus neither type of
vertex can occur. A graph not equal to its core has a
genus-zero vertex of valency at most one at the first deletion,
so this proves the claim.
\end{proof}

\begin{theorem}[Common corner refinement of core boundaries;
\status{Proved}]\label{P3a-core}
Two marked normal-crossing surface completions in the same
meromorphic class, each with boundary graph equal to its core,
have a common refinement obtained from both by corner blow-ups.
In particular, every good completion in one meromorphic class
is connected to every other by corner blow-ups and blow-downs
through good pairs.
\end{theorem}
\begin{proof}
Resolve the marked bimeromorphic map by point blow-ups supported
on the boundary, obtaining $W\to X_1,X_2$ with normal-crossing
total boundary. Resolution of a meromorphic map of smooth
surfaces and factorization of a proper birational surface
morphism into point blow-ups apply in the analytic as well as
the algebraic category; for the algebraic factorization see
\citet[Tag 0C5Q, Lemma 54.17.1]{StacksSurfaceFactorization2026}.
A point blow-up at a node subdivides an edge; at a smooth
boundary point it attaches a rational leaf. Consequently the
strict transform of every vertex of either endpoint graph lies
in the core of the graph of $W$. Every vertex outside that core
is exceptional for both projections.

If there is such a vertex, choose the one created last in a
point-blow-up factorization of $W\to X_1$. A later blow-up touching
it would create a later noncore vertex, whether at a smooth point
or at a node. Thus it has square $-1$ and valency one or two.
Contract it. The surface remains smooth and the boundary remains
normal crossing; valency two may create an allowed self-node.
Both projections descend, since they contract that curve. Repeat.
This terminates because each contraction removes one vertex.

The resulting common refinement has no fringe. In a point-blow-up
factorization from either endpoint, any blow-up at a smooth
boundary point would create a fringe which later point blow-ups
could only subdivide or enlarge, never turn into core. Hence
both factorizations consist exclusively of corner blow-ups.
For good endpoints, Theorem \ref{P2b} proves that every pair in
these two factorizations is good.
\end{proof}

\begin{theorem}[A nonemptiness test from any completion in one class;
\status{Proved}]\label{P3a-existence}
Fix any marked normal-crossing completion $(Y,B)$ of a noncompact
smooth complex surface $U$. There is a good completion in its
meromorphic class if and only if all the following tests hold.
First, $U$ is aspherical. Second, every connected component of
$B$ has nonempty core, and its entire fringe can be contracted
to smooth points, by successively contracting rational $(-1)$
tips or bivalent curves outside the core while retaining normal
crossings. Third, on the resulting core pair $(Y_c,B_c)$ the
meridian of every boundary component has infinite order in
$\pi _1(U)$ and the rank-two meridian map at every node is
injective. The contraction of the entire fringe, when possible,
is unique over its complement. Every good completion in this
class is precisely a corner modification of $(Y_c,B_c)$.
\end{theorem}
\begin{proof}
If the tests hold, all boundary curve strata of $(Y_c,B_c)$ are
aspherical: a positive-genus curve remains aspherical after
punctures are removed, and every rational curve has at least
two normalization punctures. Point strata are contractible.
Together with the first and third tests this is exactly
goodness. The common-refinement theorem then gives every good
completion in the stated class.

For necessity, resolve a bimeromorphic map from $Y$ to a good
completion $X$. On the resulting $W$, each fringe tree is
exceptional over $X$, because $X$ is core. Starting with the
last created noncore curve in a factorization over $X$ contracts
the whole fringe to smooth points and leaves a corner
refinement of $X$. Consider now a point-blow-down factorization
$W\to Y$. A core curve carrying an attached fringe cannot be
contracted while normal crossings are retained. A positive-genus
curve cannot be an exceptional $(-1)$ curve, and a rational core
vertex has core valency at least two, so an attached fringe
raises its valency to at least three. All core contractions in
this factorization are therefore disjoint from the remaining
fringe, and commute with its contraction. A noncore contraction
merely contracts part of an already smoothly contractible
exceptional tree. Induction through the factorization proves
that the fringe already visible on $Y$ is smoothly contractible.
It also shows that its contracted core model lies in the corner
class of $X$. Blow-ups preserve whether a graph component has
empty core, so no boundary component can have empty core.

Theorem \ref{P2b} transfers goodness from $X$ to this core model,
giving the first and third tests. Finally, a proper contraction
of a fixed exceptional forest to smooth points is unique: its
target structure sheaf is the direct image of the source
structure sheaf, and the resulting map is an isomorphism off
the forest. The two targets therefore identify across the
contracted points. This is an exact test from an arbitrary
starting completion, not a claim that asphericity or group
injectivity is algorithmically decidable.
\end{proof}

A \emph{minimal good completion} here means that no nonidentity
boundary contraction to another smooth normal-crossing good
completion is possible. This is boundary minimality; it is not
minimality of the underlying compact surface.

\begin{theorem}[Existence and the precise uniqueness criterion;
\status{Proved}]\label{P3a-minimal}
Every nonempty meromorphic class of good surface completions
contains a minimal good completion. Such a completion is minimal
exactly when its boundary has no smooth rational bivalent
$(-1)$ component. Let $(X,D)$ be one minimal good completion
in the class. It is the unique minimal marked completion in
that class if and only if
\[
 \begin{cases}
 C^2\leq-2 &\text{for every smooth rational bivalent component }C,\\
 C^2\leq 2 &\text{for every singly self-nodal rational bivalent
                    component }C.
 \end{cases}
\]
The second type is an isolated rational component with one
ordinary self-node; its two normalization branches count as
its valency two. Positive-genus components and rational
components with at least three punctures impose no additional
condition.
\end{theorem}
\begin{proof}
A boundary curve contractible to a smooth point is a smooth
rational $(-1)$ curve. In a core boundary it has valency at least
two, and retaining normal crossings forces valency exactly two.
Its contraction is inverse to a corner blow-up and preserves
goodness by Theorem \ref{P2b}. Repeatedly contracting such curves
terminates because the number of components decreases. This
proves both existence and the characterization of minimality.

Let $X'$ be another minimal completion and take a common corner
refinement $W$. If the marked models differ, then in a
corner-contraction sequence $W\to X'$ at least one strict
transform of an original component of $X$ disappears. Consider
the first one to disappear. Every preceding contracted curve
was exceptional over $X$, so each preceding model still maps
to $X$. A smooth original rational bivalent component therefore
has self-intersection at most its original value until it
disappears. If that value was at most $-2$, it cannot become a
$(-1)$ curve. Positive genus and valency at least three prevent
contraction for the reasons already given.

An original self-nodal rational component must first have its
node blown up before it can be contracted as a smooth curve.
Blowing up that node decreases its self-intersection by four,
since its multiplicity there is two. While the node remains
resolved and no original component has disappeared, its square
is at most $C^2-4$. If $C^2\leq2$, this is at most $-2$, again
preventing its contraction. If the resolving exceptional is
first contracted, the component becomes singular again and
still cannot be the next smooth exceptional curve. Under the
two inequalities no first original component can disappear.
Thus every contraction from $W$ to $X'$ contracts only curves
exceptional over $X$; since both endpoints are minimal, the
marked models agree.

Conversely, a smooth rational bivalent component in a minimal
pair with square $a\geq0$ can be made $(-1)$ by $a+1$ successive
corner blow-ups adjacent to it. Contract that component, then
contract any remaining bivalent $(-1)$ components until a
minimal pair is reached. Its divisorial valuation has
disappeared and cannot reappear under further contractions,
so the resulting marked minimal pair differs from $X$.
For a self-nodal rational component with square $a\geq3$,
first blow up its node. Its square is $a-4\geq-1$; perform
further adjacent corner blow-ups if necessary, then the same
contraction argument applies. This proves necessity of both
inequalities and their sufficiency.
\end{proof}

\begin{corollary}[The exact analytic scope; \status{Proved}]
\label{P3a-analytic-scope}
For an arbitrary smooth complex surface $U$, its set $\mathcal C(U)$
of good marked compactifications is the disjoint union, over
marked meromorphic completion classes admitting a normal-crossing
model, of the corner classes passing Theorem \ref{P3a-existence}.
It is empty if there is no such class. It has a unique minimal
member if and only if exactly one class passes and the inequalities
of Theorem \ref{P3a-minimal} hold there. For a fixed algebraic $U$
and algebraic markings there is only one birational completion
class, so the preceding criterion is global without this union.
If $U$ is compact, the only completion is $(U,\varnothing)$,
and it is good exactly when $U$ is aspherical.
\end{corollary}
\begin{proof}
Every boundary blow-up or blow-down is bimeromorphic and respects
the marking, so it stays in one meromorphic class. Theorem
\ref{P3a-core} and the nonemptiness and uniqueness tests give
exactly the assertions within each such class. In the algebraic
case the identity of the dense open $U$ always gives a birational
map, so there is one class. A compact open subset of a connected
Hausdorff compact surface is closed as well, and hence equals
that surface, proving the last assertion.
\end{proof}

\subsection{The six surfaces and the analytic qualifications}

\begin{lemma}[Extension of maps to a hyperbolic affine curve;
\status{Proved}]\label{P3a-Picard}
Let $D$ be a normal-crossing divisor in a complex manifold $X$,
and let $C$ be a smooth affine complex algebraic curve with
$2g(\overline C)-2+\#(\overline C\setminus C)>0$. Every holomorphic
map $X\setminus D\to C$ extends holomorphically to $X\to\overline C$.
\end{lemma}
\begin{proof}
First choose a finite etale cover $C'\to C$ whose smooth compactification
has genus at least two. If $g(\overline C)\geq2$, take the identity.
In genus zero choose three punctures as $0,1,\infty$ and use
$y^5=x(x-1)$, restricting away from all the other punctures. Its
degree-five compactification has ramification index five at those
three points, and Riemann--Hurwitz gives genus two. In genus one,
use the degree-three permutation cover with handle monodromies
$(12),(23)$, monodromy the inverse of their commutator at one
puncture, and trivial monodromies at all other punctures. The
relation in the punctured torus group is satisfied, and the action
is transitive. Compactification has one ramification point of
index three; Riemann--Hurwitz again gives genus two. Existence of
the corresponding holomorphic cover is the covering-space form
of the Riemann existence theorem.

Work locally on $(\Delta^*)^r\times\Delta^k$ near $D$. The pullback
of this finite cover has finite monodromy, so a finite coordinate
Kummer cover kills it and permits a lift of the map to $C'$.
Compose with $C'\subset\overline C'$. Give $\overline C'$ its
compact hyperbolic metric, with positive systole. For each
punctured coordinate, hold the other coordinates fixed and take
a coordinate circle of radius tending to zero. Its length in
the complete Poincare metric of $\Delta^*$ tends to zero.
Schwarz--Pick makes its image shorter than the systole for a
sufficiently small radius. Thus that generator of
$\pi _1((\Delta^*)^r\times\Delta^k)=\Z^r$ maps trivially to
$\pi _1(\overline C')$. This works for every coordinate generator.

The lifted map therefore lifts again to the universal cover
$\Delta\to\overline C'$. This holomorphic disc-valued function is
bounded, so successive applications of the removable-singularity
theorem extend it over the coordinate divisor to $\Delta^{r+k}$.
The maximum-modulus principle keeps its image in the open disc:
an attained boundary value would force constancy with modulus one,
contrary to its values on the original open set. Composing with
$\Delta\to\overline C'\to\overline C$ gives an extension on the
Kummer cover. It is invariant under the finite deck group, since
it agrees with the original pulled-back map on a dense open set;
hence it descends holomorphically. Uniqueness on overlapping
charts glues the extensions. This direct proof also supplies the
normal-crossing version of the familiar three-punctured-sphere
extension used in \citet[p.~213]{KiernanKobayashi1973}.
\end{proof}

\begin{theorem}[The six requested computations; \status{Proved}]
\label{P3a-examples}
The following describes all good marked completions, with normal
crossings allowed to include self-nodes.
\begin{enumerate}
\item For $\C^2$, for $\PP^2$ minus a smooth conic, and for
$\PP^2$ minus an irreducible nodal cubic, $\mathcal C(U)$ is empty,
in the analytic as well as the algebraic category.

\item For $U=(\C^*)^2$ with algebraic marking, the good completions
are exactly the smooth complete toric pairs. Their minimal members
are $\PP^2$ with its coordinate triangle and the Hirzebruch surfaces
$\mathbb F_n$ with toric boundary, where $n=0$ or $n\geq2$.
In the analytic category the full list consists of those toric
pairs with arbitrary holomorphic markings, and the one-curve Hopf
pairs in Ueda's classification with arbitrary holomorphic markings,
in both cases modulo marked pair isomorphism. Every one-curve Hopf
pair is already minimal and has no other good completion in its
marked meromorphic class.

\item Put $B=\PP^1\setminus\{0,1,\infty\}$ and fix the coordinates
$(x,b)$ on $U=\C^*\times B$. The minimal analytic marked completions
are exactly
\[
 X=\PP(\mathcal O\oplus L)\longrightarrow\PP^1,\qquad
 D=S_0+S_\infty+F_0+F_1+F_\infty,
\]
where $L$ is a holomorphic line bundle on $\PP^1$ equipped with a
holomorphic trivialization $\tau$ on $B$, and the marking is
$(x,b)\mapsto x\tau(b)$ in the punctured line bundle. The classes
are precisely the relative Picard group
$\operatorname{Pic}_{\rm an}(\PP^1,B)$, with
\[
 0\longrightarrow \mathcal O^*(B)/\C^*
 \longrightarrow\operatorname{Pic}_{\rm an}(\PP^1,B)
 \longrightarrow\Z\longrightarrow0.
\]
The algebraic version is
$\operatorname{Pic}_{\rm alg}(\PP^1,B)\cong
\operatorname{Div}_{\{0,1,\infty\}}(\PP^1)\cong\Z^3$.
Every good completion is a corner refinement of one of these
minimal pairs.

\item If $U=C_1\times C_2$ for two smooth affine hyperbolic curves,
every analytic marked completion belongs to the meromorphic class
of $\overline C_1\times\overline C_2$. The unique minimal good
marked completion is the product completion with its product
boundary, and every other good completion is obtained from it by
corner blow-ups. The same statement holds algebraically.
\end{enumerate}
\end{theorem}
\begin{proof}
For any noncompact surface a compactification has nonempty
boundary. Its generic normal lattice is $\Z$. Thus finite
$\pi _1(U)$ is an absolute obstruction, irrespective of the
analytic or algebraic category. The three groups in the first
item are respectively $1$, $\Z/2$, and $\Z/3$.

Here are direct calculations for the latter two. The identification
$\operatorname{Sym}^2(\PP^1)\cong\PP^2$ identifies the diagonal
with a smooth conic. The ordered configuration space of two
distinct points of $S^2$ projects to $S^2$ with contractible fibre
$S^2\setminus\{\mathrm{point}\}\cong\C$, and is simply connected.
The free interchange involution gives the unordered complement,
so its group is $\Z/2$. For a nodal cubic choose its affine form
$y^2=x^2(x+1)$. Projection to $x$ has two generic punctures in
the fibre. At $x=-1$ its braid monodromy is a half twist and
identifies their meridians; at $x=0$ it is a full twist and
imposes commutation. Van Kampen therefore gives $\Z$ for the
affine complement. Adding the line at infinity kills its
meridian, which is the product of three generic cubic meridians,
hence the third power of the generator. There is only this
additional normal relation, because the added divisor is
connected and smooth in the complement. The projective group
is $\Z/3$.

For the algebraic torus, the standard toric completion is good:
all strata are tori, and a smooth cone's primitive rays extend
to a lattice basis, so its normal meridians form an injective
coordinate sublattice of $\pi _1(U)=\Z^2$. Every algebraic marked
completion belongs to its birational class. Theorem
\ref{P3a-core} gives a common corner refinement with any good one.
Starting from the standard toric pair, these blow-ups are toric,
and their corner contractions descend the torus action.
Consequently the other pair is toric too. The minimal smooth
complete toric surface fans are the three-ray projective-plane
fan and the four-ray Hirzebruch fans, with the $(-1)$ boundary
of $\mathbb F_1$ contractible. This also follows from the explicit
cycle reduction in \citet[Proposition~3, pp.~89--90]{Ueda1979}.

For completeness we address all analytic torus completions,
rather than silently assuming they are bimeromorphic to the
standard algebraic marking. Ueda's Theorem~2 lists rational
surfaces, one-curve Hopf surfaces and their boundary blow-ups,
and the Serre ruled surface over an elliptic curve and its
boundary blow-ups \citep[p.~81]{Ueda1979}.
In the rational case $D$ is connected, since $(\C^*)^2$ has one
end. Poincare--Lefschetz duality and the homology sequence of
$(X,D)$ give a surjection
\[
 H^2(U,\Q)=H_2(X,D;\Q)\longrightarrow H_1(D,\Q),
\]
because $H_1(X,\Q)=0$. The left group has dimension one. Since
\[
 b_1(D)=2\sum_i g(\widetilde D_i)+b_1(\text{incidence graph}),
\]
all components are rational. A connected core graph with first
Betti number at most one is a cycle. Resolve a self-node and,
if needed, subdivide a two-edge cycle by corner blow-ups.
The resulting simple cycle satisfies Ueda's Proposition~3 and
reduces by corner contractions to the projective-plane or
Hirzebruch toric pair. The original pair is therefore toric
as well; its given holomorphic marking need not be algebraic.

A Hopf pair on Ueda's list has one smooth elliptic boundary.
On its primary Hopf cover the complement is $\C\times\C^*$,
and a small normal meridian is the generator of its $\Z$.
That cover is a genuine covering of the complements, so this
meridian has infinite order downstairs. The same argument
passes to secondary Hopf quotients. The open torus and the
elliptic boundary are aspherical, proving goodness. Any
boundary blow-up adds a rational fringe and is bad. Since
there are no nodes, the common-corner theorem also proves
uniqueness in its meromorphic class.

For the Serre ruled surface, the complement is an affine-line
bundle over an elliptic curve. Both that bundle and its
$\PP^1$ compactification induce the same fundamental group
as the base. Hence the normal boundary meridian dies. Boundary
blow-ups do not change the generic meridian of its strict
transform, so none can be good. This proves the complete
analytic torus list.

We next treat $\C^*\times B$. By Lemma \ref{P3a-Picard}, the
projection to $B$ extends to a holomorphic map
$p:X\to\PP^1$ on every analytic completion. Its general fibre
is the compactification of $\C^*$, hence $\PP^1$, and has
exactly two horizontal boundary points. The given coordinate
$x$ labels its two ends consistently on $B$, so these horizontal
branches have trivial permutation monodromy. Their closures
are two sections: on normalization each maps to the base
with degree one. In particular each is smooth. The Stein
factorization has degree one because the general fibre is
connected, so all fibres of $p$ are connected.

This already proves rationality without assuming a minimal
underlying surface. For one of the sections $S$, the coherent
sheaf $p_*\mathcal O_X(S)$ is torsion free on $\PP^1$, hence
locally free, and has generic rank two. Its evaluation map
defines a meromorphic map from $X$ to the projectivization of
this rank-two bundle, restricting to the degree-one complete
linear system on the general fibre. It is therefore bimeromorphic.
A holomorphic bundle over $\PP^1$ splits, so that projectivization
is a Hirzebruch surface. Resolving this bimeromorphic map and
using point factorization shows that $X$ is a smooth rational
projective surface. These auxiliary contractions are used to
prove rationality; they are not asserted to preserve the
original open $U$. This avoids an invalid appeal to a theorem
assuming the underlying compact surface was already minimal.

Consequently $p$ and $D$ are algebraic in this projective model.
Choose a rational fibre coordinate $t$ having its generic zero
and pole at the two sections. Any zero or pole of $t$ on
$U$ is vertical. The fibres of $U\to B$ are smooth irreducible
copies of $\C^*$, so
$\operatorname{div}(t)|_U$ is a finite sum of entire fibres.
Since $\operatorname{Pic}_{\rm alg}(B)=0$, multiply $t$ by a
rational base function to remove all these vertical divisors.
The resulting $t$ is an algebraic unit on $U$. On every fibre,
write it in the original coordinate as
$t=a(b)x^{k_b}$. This follows because each algebraic fibre has
compactification $\PP^1$ with two points deleted. Its winding
number $k_b$ is locally constant in $b$, and generically is
$1$ or $-1$. Thus, after inversion if needed, $k_b=1$
everywhere and $a$ is a nowhere-zero holomorphic function on
$B$. The map $(p,t):U\to B\times\C^*$ is an algebraic etale
bijection, hence an algebraic isomorphism. This supplies an
auxiliary algebraic marking; the given analytic marking
differs from it by the unit $a(b)$.

In this auxiliary algebraic marking the core theorem supplies
a common corner refinement with the product completion.
All corner operations lie over $0,1,\infty$. Neither horizontal
section can be contracted because $p$ is nonconstant on it.
In a reducible fibre arising from corner operations, the
rational chain has a bivalent $(-1)$ component until it is
irreducible; this is also the usual point-blow-up description
of a fibre of a ruled surface. Thus a minimal good model has
only smooth irreducible fibres and is a $\PP^1$ bundle with
two disjoint sections. These sections split its rank-two
bundle, giving $\PP(\mathcal O\oplus L)$. Its sections have
valency three and its three boundary fibres have square zero,
so it is boundary-minimal. Conversely each such pair is good:
over $B$ it is the marked product, the section meridians are
the nonzero $\Z$ factor, and the fibre meridians project to
nontrivial peripheral elements of $\pi _1(B)$; at their
intersections this projection makes the two meridians independent.

A marked isomorphism must preserve $p$, the zero end, and the
infinity end. It is exactly a line-bundle isomorphism carrying
$\tau$ to $\tau'$. Hence its classes are relative Picard classes.
The exact sequence displayed in the statement follows by
restricting units and line bundles from $\PP^1$ to $B$;
$\operatorname{Pic}_{\rm an}(B)=0$ because $B$ is a Stein
Riemann surface with $H^2(B,\Z)=0$. In the algebraic case a
trivialization extends to a rational section, whose divisor
is supported at the three punctures, yielding $\Z^3$.
There is no inversion quotient with the fixed marking in
the statement: $x$ labels $0$ and $\infty$. Inversion becomes
an equivalence only if one allows the marking itself to
change by $x\mapsto x^{-1}$. Holomorphic units such as
$\exp(h(b))$ explain why the analytic classification is
strictly larger than the algebraic $\Z^3$ list.

Finally, for $C_1\times C_2$, both projections extend by the
Picard lemma. Their product is a proper holomorphic map
$X\to\overline C_1\times\overline C_2$ restricting to the
specified identity on $U$, so it is bimeromorphic. The product
pair is good: curve strata are the hyperbolic affine factors,
and normal meridians are peripheral elements in independent
free-group factors. A rational boundary component has at
least three punctures; all other components have positive
genus. Thus there are no smooth or self-nodal rational
bivalent components. Theorem \ref{P3a-minimal} proves uniqueness
of the minimal model, and Theorem \ref{P3a-core} gives all its
corner refinements.
\end{proof}

\begin{proposition}[Failures of the unrestricted analytic and stack
comparison; \status{Refuted}]\label{P3a-scope}
The assertion that every two good compactifications of a smooth
complex surface are connected by stratal blow-ups and blow-downs
is false if arbitrary analytic markings are allowed. It is also
false for smooth Deligne--Mumford compactifications if the only
allowed operations are ordinary stratal blow-ups and blow-downs.
Even existence in a specified algebraic compactification class
need not agree with existence for the underlying analytic surface.
\end{proposition}
\begin{proof}
Fix $0<|\alpha|<1$, and let
\[
 H=(\C^2\setminus\{0\})/
      \langle(z,w)\mapsto(\alpha z+w,\alpha w)\rangle,\qquad
 E=\{w=0\}/\langle g\rangle.
\]
On $w\ne0$, put $u=\alpha z/w$ and $v=w$. Then
$g(u,v)=(u+1,\alpha v)$. Choosing a logarithm of $\alpha$,
the invariant coordinates
\[
 A=e^{2\pi i u},\qquad B=v\,e^{-u\log\alpha}
\]
identify $(H\setminus E)$ biholomorphically with $(\C^*)^2$.
The pair is good by the Hopf argument above. It has $b_1(H)=1$,
whereas a toric rational completion has $b_1=0$.
Blowing up points preserves $b_1$, so these two good analytic
pairs cannot be connected by the requested operations.
Their dual complexes are respectively a point and a circle.
Even on a fixed toric surface, changing the marking by
$(x,y)\mapsto(x,e^x y)$ need not give the same meromorphic
class: an alleged meromorphic extension would restrict on a
generic $y$-section to the meromorphic continuation of $e^x$
across $x=\infty$, which is impossible.

For the stack assertion, start with a good toric pair and
take a nontrivial root of a boundary divisor. Its strata are
finite gerbes over torus strata and are aspherical in the
stack sense. On an atlas its meridian is a nonzero multiple
of the old primitive meridian, so the pair is good.
Triviality of all inertia is preserved by an ordinary
representable blow-up in both directions. One direction follows
because blowing up an algebraic space again has trivial inertia.
For the converse, if a nonidentity finite stabilizer element $g$
fixes a point outside the centre, it persists upstairs. If the
point lies in the centre, the blow-up fibre is the projective
space of its normal representation. Over $\C$, $g$ has an
eigenline, so some point of that fibre still has nontrivial
stabilizer. Thus no zigzag of these blow-ups can connect the
root pair to the toric pair with trivial inertia. Generic
inertia on a particular boundary divisor is not claimed to be
invariant: exceptional divisors can acquire such inertia.

For the last assertion let
$H_q=(\C^2\setminus\{0\})/\langle(z,w)\mapsto(qz,qw)\rangle$,
$0<|q|<1$, and delete $E_1=\{w=0\}/q^\Z$.
Then
\[
 H_q\setminus E_1\cong E_q\times\C,\qquad
 E_q=\C^*/q^\Z,\qquad
 (z,w)\longmapsto([w],z/w).
\]
The boundary $E_1$ is elliptic. Near it, use $[z]$ and
$\xi=w/z$; its normal meridian maps to
$([\xi z],1/\xi)$. The angular loop in $E_q$ is nonzero
and of infinite order, so this is a good analytic completion.
For the fixed algebraic surface $E_q\times\mathbb A^1$, no
algebraic good completion exists. In any proper algebraic
completion, the map $\mathbb A^1\to U$ given by one affine
fibre extends to $\PP^1\to X$. The sole missing point has
positive contact order with at least one boundary branch.
A sufficiently large circle bounds the finite affine disc
inside $U$, but maps to a nonzero full normal-meridian
vector at that missing point. This contradicts goodness.
Equivalently the standard algebraic completion has a
boundary meridian which dies, and the core test forbids all
other algebraic completions. The Hopf and algebraic ruled
surfaces are not bimeromorphic: their first Betti numbers
are one and two. Thus a meromorphic-class qualification
is indispensable even for nonemptiness.
\end{proof}

\subsection{Elementary fibrations, stable boundaries, and arithmetic inertia}

An elementary hyperbolic fibration here means a smooth family of affine
curves obtained from a smooth proper family by removing a nonempty
finite etale marked divisor, with $2g-2+n>0$. Analytically it is a
locally trivial bundle with fibre a punctured compact surface; its
fundamental group is a finitely generated nonabelian free group.
An iterated tower starts from a point. Base points and paths are
specified when homomorphisms, rather than their conjugacy classes,
are wanted.

\begin{theorem}[Horizontal inertia over an arbitrary base;
\status{Proved}]\label{P3b-discrete}
For an elementary hyperbolic fibration $F\to U\to B$ with connected
base, the fibre map $\pi _1(F)\to\pi _1(U)$ is injective. Every
horizontal puncture meridian has infinite order in $\pi _1(U)$.
Neither assertion requires $B$ to be aspherical. An iterated
hyperbolic tower is aspherical.
\end{theorem}
\begin{proof}
The possible obstruction to fibre injectivity is the connecting map
$\partial:\pi _2(B)\to\pi _1(F)$. Pull the bundle back along a
representative sphere. Its clutching loop is a loop $h_t$ of
self-homeomorphisms of $F$, based at the identity. At a chosen point
$x$, the loop $h_t(x)$ represents $\partial$. For any based loop
$\gamma(s)$ in $F$, the square $h_t(\gamma(s))$ supplies a homotopy
showing that $h_t(x)$ commutes with $\gamma$. Hence the image of
$\partial$ is central. The centre of a nonabelian free group is
trivial: a reduced word commuting with a free generator is its
power, and the cyclic subgroups of two different generators
intersect trivially. Thus $\partial=0$.

A puncture circle in a hyperbolic affine curve is nontrivial.
For at least two punctures it survives in homology; for one
puncture and positive genus it is the nontrivial reduced product
of handle commutators in the free group. A free group is
torsion-free, so every such meridian is infinite. The fibre
injection proves the horizontal assertion. Finally, the homotopy
exact sequence of a bundle with aspherical base and fibre gives
asphericity of its total space. Induction proves the tower
assertion.
\end{proof}

\begin{lemma}[Free normal subgroups and profinite completion;
\status{Proved}]\label{P3b-completion}
If $1\to N\to G\to Q\to1$ is an extension and $N$ is a finitely
generated nonabelian free group, then
\[
 1\longrightarrow\widehat N\longrightarrow\widehat G
 \longrightarrow\widehat Q\longrightarrow1
\]
is exact. Every nontrivial cyclic subgroup of $N$ induces an
embedding $\widehat\Z\to\widehat N$.
\end{lemma}
\begin{proof}
We first prove $Z(\widehat F_r)=1$ for $r\ge2$, rather than
assuming that discrete centre-freeness survives completion.
Let $F_r\to A$ be any finite quotient with kernel $K$, and choose
a prime $p\nmid |A|$. The regular covering of a rose with $r$
edges has cellular exact sequence
\[
 0\longrightarrow H_1(K;\mathbf F_p)
 \longrightarrow\mathbf F_p[A]^r
 \longrightarrow\mathbf F_p[A]
 \longrightarrow\mathbf F_p\longrightarrow0.
\]
Averaging over $A$ splits every sequence of these modules.
Consequently
\[
 H_1(K;\mathbf F_p)\cong
 \mathbf F_p\oplus\mathbf F_p[A]^{r-1},
\]
on which $A$ acts faithfully. The finite group
$F_r/[K,K]K^p$ is an extension of $A$ by this module, with this
conjugation action. The image of a central element of
$\widehat F_r$ in that finite group acts trivially on its kernel;
its image in $A$ is therefore $1$. Since $A$ was arbitrary,
the central element is $1$.

For a topologically finitely generated profinite group, its
continuous automorphism group is profinite. Indeed, intersections
of all open subgroups of index at most $n$ are characteristic
open subgroups forming a basis, and compatible automorphisms of
the resulting finite quotients form a closed inverse-limit group.
Conjugation gives
\[
 G\longrightarrow\operatorname{Aut}(N)
 \longrightarrow\operatorname{Aut}(\widehat N),
\]
which extends to $\widehat G$. Its composite with
$\widehat N\to\widehat G$ is the inner-automorphism map of
$\widehat N$. The centre-freeness just proved makes this
composite injective, proving $\widehat N\hookrightarrow\widehat G$.
The image is the closure of $N$, and quotienting $\widehat G$
by that closure gives $\widehat Q$, since both classify precisely
the finite quotients of $G$ killing $N$. This proves exactness
without a goodness hypothesis on $Q$.

For the cyclic assertion, represent a nontrivial conjugate of
$w\in F_r$ by a cyclically reduced word and its immersed labelled
circle in the rose. Complete that finite immersion to a finite
cover: for each edge label, pair the finitely many missing
incoming and outgoing edges. Their numbers agree. The original
circle is then an embedded cycle in a finite covering graph.
Its homology class is primitive, so the covering subgroup $H$
admits a homomorphism $H\to\Z$ sending $w$ to $1$.
For any $n$, take the kernel of $H\to\Z/n$ and then its core
in $F_r$. The resulting finite quotient of $F_r$ gives $w$
order divisible by $n$. Thus the topology induced on
$\langle w\rangle$ contains every congruence topology $n\Z$,
which is exactly the injectivity of
$\widehat\Z\to\widehat F_r$.
\end{proof}

\begin{corollary}[Horizontal profinite injection and towers;
\status{Proved}]\label{P3b-horizontal-profinite}
In Theorem~\ref{P3b-discrete}, every horizontal meridian induces
an injection $\widehat\Z\hookrightarrow\widehat{\pi _1(U)}$,
over any connected base. The group of an iterated elementary
hyperbolic tower is residually finite.
\end{corollary}
\begin{proof}
The homotopy sequence and the preceding theorem give a group
extension with kernel $\pi _1(F)$, so apply
Lemma~\ref{P3b-completion} twice, first to the fibre and then to
its nontrivial peripheral cyclic subgroup. For residual
finiteness, induct on the tower: an element nontrivial in the
base is separated by a finite quotient of the base, and a
nontrivial element in the fibre survives in its embedded
profinite completion. The latter group is residually finite
by the same finite-cover construction. This proves the induction.
The related goodness and completion statements for elementary
anabelian towers appear in
\citet[Lemmas 10.9 and 10.11]{Bresciani2021}; here the affine
free-fibre argument is supplied explicitly, and the horizontal
statement has the stronger arbitrary-base scope.
\end{proof}

\begin{theorem}[The corrected vertical statement;
\status{Proved}]\label{P3b-vertical}
Suppose the fibration extends to a compatible stable pointed
family over an NC base completion, after the indicated finite
root or quotient base changes, and take a smooth toroidal model
of the total boundary. Assume:
\begin{enumerate}
\item the base's open strata are aspherical;
\item its discrete full meridian lattices inject in $\pi _1(B)$;
\item for the profinite assertion, its completed full meridian
lattices also inject in $\widehat{\pi _1(B)}$.
\end{enumerate}
Then the total completion is good, and under the third hypothesis
its full completed meridian lattices inject as well. These
hypotheses hold inductively for recursively compatible
stable/toroidal completions of elementary hyperbolic towers
whose initial curve completion has its usual peripheral
injections. Stable extension alone does not imply either
vertical conclusion for an arbitrary chosen base completion.
\end{theorem}
\begin{proof}
Stable curve components, punctured at their nodes and markings,
are aspherical. A rational component has at least three special
points; an elliptic component has a special point unless it is
a whole unpointed genus-one curve, excluded from the hyperbolic
stable pointed problem. Toroidal resolution can additionally
insert rational bridges with two punctures, which are
$\C^*$. Compatibility includes constancy of the fibre incidence
type on each base stratum. The total open strata are therefore
bundles of these curves, points, or tori over finite incidence
covers of base strata. Their homotopy sequences prove their
asphericity under the first hypothesis; the open total space is
aspherical when the open base is, by Theorem~\ref{P3b-discrete}.

Consider a local total normal lattice $A$ and the corresponding
base lattice $A_B$. Projection induces an integral map $M:A\to A_B$.
At a smooth vertical point it has zero kernel; at a horizontal
marking its kernel is the fibre's puncture circle. At a node
the basic local equation is $xy=t$, for which
\[
 \Z^2\longrightarrow\Z,\qquad (u,v)\longmapsto u+v,
\]
and the kernel generator $(1,-1)$ is the vanishing cycle.
After finite base change the equations are monomial; smooth
toroidal charts replace the displayed lattices by primitive
sub-lattices and finite-index lattices. Thus at every point
the kernel has rank at most one and, if present, maps to a
nonzero power of a puncture or vanishing cycle in $\pi _1(F)$.

The relevant vanishing cycle is nontrivial. A smoothing neck
could bound a disc with no puncture only if the special fibre
had an unstable unmarked rational tail on that side; stable
reduction removes it. A peripheral neck enclosing one puncture
is still nontrivial, as are the necks in inserted two-punctured
bridges. All other necks are essential surface curves.
This verifies the kernel claim also in the separating case.

Now use the commutative diagram with exact rows
\[
\begin{array}{ccccc}
 \ker M&\longrightarrow&A&\longrightarrow&A_B\\
 \downarrow&&\downarrow&&\downarrow\\
 \pi _1(F)&\longrightarrow&\pi _1(U)&\longrightarrow&\pi _1(B).
\end{array}
\]
The right vertical arrow is injective by hypothesis, and the
left arrow is injective because a nontrivial free-group element
has infinite order. Any element killed by the middle arrow
first lies in $\ker M$ and then is zero. This proves the full
discrete lattice assertion.

Complete the bottom row using Lemma~\ref{P3b-completion}.
The kernel of the completed integral matrix $M$ is the
completion of its integral kernel: the image is a free abelian
group, the short exact sequence with that image splits, and
its inclusion in $A_B$ is still injective after completion.
The third hypothesis and the cyclic part of the same lemma
now give the identical diagram chase with $\widehat A$.
Starting from a punctured hyperbolic curve, these discrete and
completed peripheral injections hold by the finite-cover
argument above. Roots restrict to finite-index lattices, and
stratal subdivisions use unimodular charts, so the induction
survives all the specified compatible operations. The theorem
is conditional on such a compatible stable/toroidal completion;
it does not silently assert that every arbitrarily chosen
completion of a tower is compatible.
\end{proof}

\begin{example}[Stable extension on a bad base; \status{Refuted}]
\label{P3b-vertical-counterexample}
Take $B=(\PP^1\setminus\{0,1,\infty\})^2$ with its product
completion, and blow up a smooth point of one boundary component.
At the new boundary node the strict-transform and exceptional
meridians both map to the same old meridian. Their difference
is killed. For any hyperbolic affine curve $F$, the constant
stable marked family $\overline F\times Y\to Y$ over this
blown-up base is already stable; its open total space is the
genuine three-stage hyperbolic tower $F\times B$.
Nevertheless the same rank-two vertical kernel persists in
$\pi _1(F)\times\pi _1(B)$, and in its profinite completion.
This refutes the unrestricted vertical assertion even within
the requested class of towers. It also illustrates why
injectivity of individual meridians is inadequate.
\end{example}

\begin{corollary}[Characteristic-zero geometric inertia;
\status{Proved}]\label{P3b-profinite}
Let the data be defined over a field $K$ finitely generated over
$\Q$, and choose $\overline K\hookrightarrow\C$. Under the
hypotheses just established, horizontal geometric inertia,
and the full vertical inertia in the compatible tower case,
inject into $\pi _1^{\mathrm{\acute et}}(U_{\overline K})$:
\[
 T^{I(x)}=\widehat\Z(1)^{I(x)}
       \hookrightarrow\pi _1^{\mathrm{\acute et}}(U_{\overline K}).
\]
These are homomorphisms up to the conjugations determined by
geometric base-point paths; no preferred trivialization of
the Tate module is asserted.

If the local punctured-neighbourhood extension
$1\to I_x\to E_x\to G_{K(x)}\to1$ maps to
$1\to\pi _1(U_{\overline K})\to\pi _1(U)\to G_K\to1$
and the chosen $G_{K(x)}\to G_K$ is injective, then
\[
 \ker(E_x\to\pi _1(U))
 =\ker(I_x\to\pi _1(U_{\overline K})).
\]
Thus in that case full local arithmetic injectivity is
equivalent to geometric inertia injectivity. This qualification
holds for a closed point after the corresponding finite field
extension, with compatible separable closures. It must not
be omitted for a transcendental residue field.
\end{corollary}
\begin{proof}
Riemann existence identifies the geometric etale fundamental
group with the profinite completion of the complex topological
one; the precise comparison is
\citet[Exposé XII, Théorème 5.1 and Corollaire 5.2]{SGA1}.
In an NC chart the compatible Kummer covers of the punctured
coordinates give $T^{I(x)}$. Under the chosen complex embedding
and a compatible system of roots of unity, this is precisely
the completed normal lattice used above. The proved profinite
injections therefore give the assertion. Changing paths
conjugates the maps; Galois acts cyclotomically and may permute
branches, which is why the intrinsic Tate-module notation is
retained.

For the final statement, project any element of
$\ker(E_x\to\pi _1(U))$ to $G_{K(x)}$. Its image in $G_K$ is
trivial, hence its image in $G_{K(x)}$ is trivial by the stated
hypothesis. It belongs to $I_x$, where the two kernels agree.
The reverse inclusion is immediate. For a transcendental
residue field the Galois-quotient map need not be injective,
so this argument gives no such equivalence. None of these
characteristic-zero arguments treats wild inertia in positive
characteristic.
\end{proof}

\subsection{Higher-dimensional comparison: established results and the remaining question}

\begin{proposition}[An additional projective obstruction;
\status{Proved}]\label{P3c-Mori}
A good smooth projective NC pair admits no nonconstant algebraic
map $\mathbb A^1\to S$ into any of its open strata, including $U$.
For a simple normal-crossing boundary it follows that
$K_X+D$ is nef.
\end{proposition}
\begin{proof}
Let $f:\mathbb A^1\to S$ be nonconstant. Properness extends it to
$\overline f:\PP^1\to X$, with image in the closure of $S$.
If the image of infinity remains in $S$, asphericity makes this
sphere nullhomotopic in $S$, hence in $X$. But the pullback of
a Kähler form of the projective $X$ has strictly positive
integral on a nonconstant holomorphic sphere, a contradiction.

Otherwise infinity maps to a deeper stratum $T$. Pull back the
normal directions of $S$ along $\mathbb A^1$. Their bundles and branch
labelling are topologically trivial because $\mathbb A^1$ is contractible.
Choose nonzero normal displacements and a continuous positive
radius that tends to zero faster than every newly vanishing
boundary coordinate near infinity. This pushes $f(\mathbb A^1)$ into
$U$. Its large circle bounds the pushed finite disc. In a local
NC chart at $\overline f(\infty)$, however, the circle has a
nonzero meridian vector: on every branch newly acquired at
infinity the winding is the nonzero order of vanishing, with
a common orientation sign. The radius choice preserves these
windings. The coordinates in the old normal directions can be
arbitrary integers, but cannot cancel the distinct new
coordinates. Thus a nonzero element of the full lattice at
$T$ dies in $\pi _1(U)$, contrary to goodness. This proves
the first assertion, also when a normalization of the stratum
closure is needed to keep track of branches.

For a reduced SNC divisor these open strata are exactly those
in the Mori-hyperbolicity hypothesis of
\citet[Theorem 1.1]{Svaldi2019}. That theorem, applied to the
smooth projective log canonical pair $(X,D)$, gives
$K_X+D$ nef. This last step uses the cone theorem; no claim
that asphericity alone implies positivity is being made.
\end{proof}

\begin{theorem}[Birational morphisms between good projective pairs;
\status{Proved}]\label{P3c-morphism}
Let $f:(Y,E)\to(X,D)$ be a proper birational morphism of smooth
projective varieties with reduced SNC boundaries, inducing the
specified identity on $U=Y\setminus E=X\setminus D$. If both pairs
are good, then $f$ is logarithmically etale and toroidal. Moreover
\[
 K_Y+E=f^*(K_X+D).
\]
This asserts toroidality of the morphism, not a factorization of
it solely by smooth stratal blow-ups.
\end{theorem}
\begin{proof}
The natural logarithmic differential map
\[
 f^*\Omega_X^1(\log D)\longrightarrow\Omega_Y^1(\log E)
\]
is defined everywhere. Indeed, in boundary coordinates each
$f^*x_i$ is a unit times a monomial in the local equations of
$E$, so its logarithmic differential has only logarithmic poles.
It is an isomorphism on $U$. Its determinant has an effective
divisor $A$, and
\[
 K_Y+E=f^*(K_X+D)+A.
\]
At the generic point of every nonexceptional prime divisor a
proper birational morphism between smooth varieties is an
isomorphism. Thus $A$ is $f$-exceptional. Proposition
\ref{P3c-Mori} makes $K_Y+E$ nef, so $A$ is $f$-nef. The
negativity lemma applied to $-A$ gives $-A\geq0$, because
$f_*(-A)=0$. With $A\geq0$ this gives $A=0$. We use the sign
convention of \citet[Lemma~3.39]{KollarMori1998}, also stated
in \citet[Lemma~1.73, author version p.~51]{Xu2025}.
The determinant has no zero, so the logarithmic differential
map is an isomorphism everywhere.

Here is the local chart implication, in addition to the abstract
criterion \citep[Proposition~3.12 and Theorem~3.5]{Kato1989}.
Let $x_1,\ldots,x_r$ be the target boundary coordinates and
$y_1,\ldots,y_s$ the source boundary coordinates at the points
in question. Write
\[
 f^*x_i=u_i\prod_{j=1}^s y_j^{a_{ij}},
 \qquad i=1,\ldots,r,
\]
with units $u_i$, and write $g_\ell=f^*x_\ell$ for the remaining
ordinary coordinates. In logarithmic bases, at the source
point the derivative has block form
$\left(\begin{smallmatrix}A&B\\0&C\end{smallmatrix}\right)$,
where $A=(a_{ij})$ has $s$ columns. Invertibility forces
$\operatorname{rank}A=s\leq r$. Choose $s$ pivot rows $I$
whose matrix $A_I$ has nonzero determinant. Locally analytically
take logarithms of their units and put
\[
 z_j=y_j\exp\bigl((A_I^{-1}\log u_I)_j\bigr).
\]
Then the pivot target coordinates are pure monomials in the
$z_j$; every other boundary coordinate is a monomial times a
unit $h_i$. Row elimination in the same logarithmic Jacobian
shows that
\[
 z_1,\ldots,z_s,\quad h_i\ (i\notin I),\quad
 g_{r+1},\ldots,g_n
\]
are ordinary local coordinates, with the $h_i$ regarded as
invertible coordinates near their nonzero values. Algebraically
one obtains the same change after adjoining finitely many
roots of units, which is etale in characteristic zero.

Consequently the chart is induced by
$\mathbb N^r\to\mathbb N^s\oplus\Z^{r-s}$, with the nonpivot
units supplying the last generators. Its group matrix has
determinant $\pm\det A_I$, hence finite kernel and cokernel
of invertible order, and the map to the associated monoid
base change is ordinarily etale in the displayed coordinates.
Kato's chart criterion gives logarithmic etaleness. These
are exactly etale-local toroidal monomial charts, establishing
toroidality rather than merely vanishing relative differentials.
\end{proof}

\begin{proposition}[Crepancy for rational comparisons;
\status{Proved}]\label{P3c-crepant}
Any two smooth projective good SNC completions of the same
algebraically marked $U$ are log-crepant: on a common smooth
NC resolution $p:W\to X$, $q:W\to Y$,
\[
 p^*(K_X+D)=q^*(K_Y+E).
\]
The effective logarithmic Jacobian divisors of the two
projections are equal and are exceptional for both.
\end{proposition}
\begin{proof}
Let $B=W\setminus U$ with reduced structure. As in the preceding
proof, the two logarithmic Jacobians give
\[
 K_W+B=p^*(K_X+D)+A_p=q^*(K_Y+E)+A_q,
\]
where $A_p,A_q$ are effective and exceptional for their respective
projections. Put
$R=p^*(K_X+D)-q^*(K_Y+E)=A_q-A_p$.
Both endpoint log canonical divisors are nef by Proposition
\ref{P3c-Mori}. Thus $-R$ is $p$-nef, and
$p_*R=p_*A_q$ is effective. The negativity lemma gives $R\geq0$.
Likewise $R$ is $q$-nef and $q_*(-R)=q_*A_p$ is effective, so
the same lemma gives $-R\geq0$. Hence $R=0$ and $A_p=A_q$.
Equality does not force this common effective divisor to vanish:
$W$ has not been assumed good. Contracting all of its positive
support compatibly with both projections is precisely a gap
which these arguments do not fill.
\end{proof}

\begin{theorem}[The complete toric comparison; \status{Proved}]
\label{P3c-toric}
Let $X_\Sigma$ and $X_{\Sigma'}$ be smooth complete toric
varieties with their full toric boundaries, marked by the
same open algebraic torus $T$. In every dimension, and in
particular for threefolds, the pairs are good and are joined
over $T$ by a finite zigzag of smooth stratal blow-ups and
blow-downs through good pairs.
\end{theorem}
\begin{proof}
Each open stratum is a torus, hence aspherical. If a smooth
cone is generated by primitive rays $v_1,\ldots,v_r$ in the
cocharacter lattice $N$, these vectors extend to an integral
basis. The local normal-meridian map is the inclusion
\[
 \Z^r\longrightarrow N=\pi _1(T),\qquad e_i\longmapsto v_i,
\]
so it is injective, both discretely and after completion.

A blow-up of the closure of the orbit associated with the
cone $\langle v_1,\ldots,v_r\rangle$ replaces it by the smooth
star subdivision with new ray $v_1+\cdots+v_r$. Thus ordinary
smooth stratal blow-ups are precisely these smooth star
subdivisions; their inverses are smooth star assemblies.
The toric weak factorization theorem supplies a finite
sequence of these subdivisions and assemblies between two
nonsingular fans with the same support. The precise theorem
used is \citet[Theorem 4.7.1]{Wlodarczyk2009}; its hypotheses
hold because both fans are nonsingular and complete in the
same $N_\R$. Every intermediate fan has that support, hence
is complete, and the preceding stratum and lattice calculation
makes every intermediate pair good. This proves the entire
requested toric case, not just preservation under a single
toric blow-up.
\end{proof}

\begin{proposition}[The weak and strong requests are different;
\status{Proved}]\label{P3c-weak-strong}
The toric zigzag theorem and the weighted strong factorization
theorem do not establish a common roof obtainable solely by
smooth stratal blow-ups.
\end{proposition}
\begin{proof}
The preceding zigzag permits blow-ups and blow-downs in
alternating order. Requiring a single common refinement
reachable from both fans solely by smooth star subdivisions
is the unweighted strong factorization problem. It is the
explicit question in
\citet[arXiv version 2, Section 7.1, Question 7.1]{AdiprasitoPak2026}.
The weighted theorem now published by those authors allows
arbitrary rational stellar subdivisions. Such a subdivision
can introduce a ray different from the sum of primitive
generators, and can pass through singular toric varieties;
it does not prove the common smooth roof assertion. The
surface common-roof theorem above therefore must not be
promoted to higher dimensions by confusing these two results.
\end{proof}

\begin{conjecture}[A restricted comparison question;
\status{Open}]\label{P3c-general}
Let $U$ be a fixed smooth algebraic threefold over $\C$, and
let $(X,D)$ and $(Y,E)$ be smooth proper good NC completions
with algebraic markings by $U$. Are they joined over $U$ by
a finite zigzag of blow-ups and blow-downs of smooth closures
of boundary strata, all intermediate pairs being good?
\end{conjecture}

The unrestricted analytic and DM-stack versions are refuted
by Proposition~\ref{P3a-scope}; even the analytic surface
version is false. The displayed scheme question is a
deliberately restricted replacement. We have proved it for
surfaces and for toric pairs. Ordinary weak factorization,
even with centres having normal crossings with the boundary,
does not answer it: Theorem~\ref{P2a} shows exactly why its
non-stratal centres are forbidden here. Propositions
\ref{P3c-Mori} and \ref{P3c-crepant} give nefness and log-crepancy,
and Theorem \ref{P3c-morphism} proves toroidality for an actual
birational morphism. They do not provide a common good resolution
of an arbitrary rational comparison. We give no proof or
counterexample for the remaining nontoric threefold question
and do not assert that its formulation constitutes a theorem
or a proof of its status in the literature.

\begin{proposition}[A failed elementary flop candidate;
\status{Proved}]\label{P3c-flop-test}
Consider an ordinary Atiyah flop between smooth projective pairs,
with exceptional curve
$C\cong\PP^1$, normal bundle $\mathcal O(-1)^{\oplus2}$.
Assume locally that exactly one boundary surface contains
$C$, that $C$ has square $-1$ on that surface, and that the
other boundary branches meet $C$ transversely at $n$ distinct
points, with no additional boundary near $C$. This construction
cannot give a good-to-good NC comparison with a non-stratal
forward blow-up.
\end{proposition}
\begin{proof}
For $n=0$, the containing open stratum contains the complete
rational curve, contrary to asphericity and Kähler positivity.
For $n=1$, its punctured part is an affine line, contrary to
Proposition~\ref{P3c-Mori}. In the common blow-up, the exceptional
quadric is $\PP^1\times\PP^1$. The containing boundary surface
meets it in a ruling fibre, whereas each transverse boundary
branch meets it in a fibre of the other ruling. After the
flop, each of those $n$ branches contains the new curve $C'$.
For $n\ge3$ this violates normal crossings along a codimension-two
curve. For $n=2$, the inverse blow-up centre $C'$ is the closure
of the intersection stratum of the two branches, hence is
stratal. If the target pair were good, its blow-up would
therefore be good by Theorem~\ref{P2b}. But the forward
blow-up has centre properly inside a single boundary surface
and is non-stratal, so is bad by Theorem~\ref{P2a}. The common
blow-up cannot be both. This excludes this precise candidate,
without treating it as a classification of all threefold
flops or a solution of Conjecture~\ref{P3c-general}.
\end{proof}

\section{Boundary maquettes, dual complexes, and tropical links}
\label{sec:boundary}

Write $D=X\setminus U$ and $L=U_\infty=\partial X^{\log}$, where the boundary of the oriented real blow-up is rounded when an ordinary manifold boundary is needed. The flag and homotopy-colimit constructions below apply to complex manifolds and smooth analytic Deligne--Mumford stacks with the étale presentations fixed in the conventions. Their topological realizations retain all discrete inertia, including infinite and ineffective inertia. The compact orientation argument explicitly assumes a proper finite-isotropy orbifold; the weight statements explicitly assume proper algebraicity. Here \emph{coarse} describes the ordinary branch-incidence realization and does not assert existence of a coarse analytic moduli space.

Let $\calD_d$ be the closed normalized incidence space of an unordered set of $d$ distinct local branches of $D$. Locally it is the disjoint union of their $d$-fold intersections. For $d_0>\cdots>d_q>0$, let $\calD_{d_0\cdots d_q}$ parametrize flags $I_0\supset\cdots\supset I_q$ of chosen branch sets, with $|I_j|=d_j$. These constructions descend étale locally, including their stabilizers. The positive-rank restriction is essential: adjoining the empty branch set adds an ambient cone vertex.

\begin{proposition}[Components do not preserve the flag-category Segal identities]
\label{P4a1}
\status{Refuted}. Taking connected components of the object and arrow spaces of the internal incidence category need not produce a category. In particular, its full flag spaces cannot in general be recovered from successive pairs of component arrows.
\end{proposition}
\begin{proof}
Let $\tau=(12)(34)$ act on $\C^4\times\C^*$ by
\[
 (z_1,z_2,z_3,z_4,w)\longmapsto(z_2,z_1,z_4,z_3,-w).
\]
The action is free. Thus its quotient is a smooth manifold, and the coordinate hyperplanes descend to a normal-crossings divisor. The connected components of its normalized flag spaces are exactly the $\langle\tau\rangle$-orbits of branch flags: each coordinate intersection times $\C^*$ is connected.

Use the rank-two set $B=\{1,2\}$. The upper pairs $(123,B)$ and $(124,B)$ belong to one component of $\calD_{32}$; the lower pairs $(B,1)$ and $(B,2)$ belong to one component of $\calD_{21}$. Their images in $\pi_0\calD_2$ agree. Nevertheless,
\[
 (123,B,1),\qquad (123,B,2)
\]
represent distinct components of $\calD_{321}$: applying $\tau$ changes $123$ to $124$, and hence cannot identify these representatives. Their composites $(123,1)$ and $(123,2)$ also lie in different components of $\calD_{31}$. Consequently the map
\[
 \pi_0\calD_{321}\longrightarrow
 \pi_0\calD_{32}\times_{\pi_0\calD_2}\pi_0\calD_{21}
\]
has a two-element fibre over the selected pair, and composition would have two different values. The internal analytic fibre-product identity remains correct; the invalid operation is commuting $\pi_0$ with that fibre product. This distinction applies to the category described in MPL~147, \S11, fol.~56, scan~114. The program \texttt{boundary\_checks.py} enumerates these orbits and their distinct composites.
\end{proof}

\begin{theorem}[The corrected component realization and its homotopy fibres]
\label{P4a2}
\status{Proved}. Define a semisimplicial set by
\[
 K_q=\coprod_{d_0>\cdots>d_q>0}\pi_0\calD_{d_0\cdots d_q},
\]
with face maps forgetting entries of the flag. Its realization is the barycentric subdivision of the ordinary coarse dual complex $\Delta_D$. Let $A_K$ be its category of simplices and specified face inclusions, with arrows from a simplex to its faces. There are diagrams of homotopy types $Y_D,Y_L:A_K\to\mathrm{Spaces}$ and natural equivalences
\[
 \hocolim_{A_K}Y_D\simeq D,\qquad
 \hocolim_{A_K}Y_L\simeq L,\qquad BA_K\simeq\Delta_D.
\]
Here $Y_D(a)$ is the connected normalized flag component represented by $a$, and $Y_L(a)$ is its normalized closed oriented face. The latter is formed by fibre product with $L$ in topological stacks before passing to homotopy types. Projection to the indexing category gives maps
\[
 q_D:D\longrightarrow\Delta_D,
 \qquad q_L:L\longrightarrow\Delta_D
\]
compatible with the blow-down $L\to D$.

Put $G_D=\Pi_\infty\Delta_D$. There are functors
\[
 \Exit(D)\longrightarrow G_D,\qquad \mathcal H_X\longrightarrow G_D,
\]
where $\mathcal H_X$ is the crossing $\infty$-category with classifying space $L$, or its ordinary model when the criterion of Part~1 applies. For either $Y=Y_D,Y_L$ and $b\in G_D$, set
\[
 A_{K,b}=A_K\times_{G_D}(G_D)_{/b}.
\]
Then the exact fibre formula is
\begin{equation}
 \label{eq:boundary-fibre}
 \hofib_b(q_Y)\simeq
 \operatorname*{colim}_{(a,\gamma)\in A_{K,b}}Y(a).
\end{equation}
The indexing pullback is taken in $\infty$-categories; $\gamma$ is a path from the image of $a$ to $b$, with all its coherences. For a functor $F:\mathcal E\to G_D$, its categorical homotopy fibre is
$\mathcal E_b=\mathcal E\times_{G_D}(G_D)_{/b}$, and
\[
 B\mathcal E_b\simeq\hofib_b(BF).
\]
The categorical fibre itself should not be identified with its groupoid completion or with a single geometric stratum.
\end{theorem}
\begin{proof}
In a chart with branch set $I$, the dual simplex has vertices $I$. A simplex of its barycentric subdivision is a strict chain of nonempty subsets of $I$. Transition maps permute branches. An automorphism preserving such a chain fixes every barycentric vertex and hence its whole barycentric simplex pointwise. Therefore taking ordinary quotients after this subdivision gives precisely the semisimplicial set of components of full flags. This proves the assertion about $K$, including self-intersections and branch monodromy. The nerve of $A_K$ gives one further subdivision. Strict flags give a semisimplicial object; a full simplicial nerve requires repeated ranks and the corresponding identities.

For the first homotopy-colimit assertion, the local augmentation is the closed-cover resolution by coordinate branch intersections, with its flag subdivision. Coordinate intersections and their inclusions are cofibrant closed-face data: collars, or a compatible triangulation of the coordinate arrangement, give the usual homotopy-colimit nerve lemma. Equivalently, over a point the incidence poset is the nonempty subsets of its local branches and has the full branch set as an extremal object; the collar models make this contraction a local equivalence of the augmented resolutions. Gluing these local equivalences by an open Čech cover proves the assertion for manifolds. For $L$, the same construction uses closed faces of the orthant of radial coordinates, retaining the angular factors. Their inclusions have collars, and their union is the boundary orthant. The same nerve lemma proves the second assertion. Forgetting the angular coordinates is a natural transformation of these diagrams and supplies the compatibility of $q_L$ and $q_D$.

Apply the local manifold argument on every level of an étale Čech nerve. Normalized branch flags and closed oriented faces commute with étale pullback, and the two resulting homotopy colimits commute. Descent therefore yields the asserted equivalences for the stack itself. This uses the full Čech nerve, retains all inertia, and imposes no finiteness condition on stabilizers. Branch permutations are nevertheless finite, since they act on a finite local branch set.

The functors in the statement are the composites with groupoid completion and the maps $q_D,q_L$. To compute a fibre, regard $A_K\to G_D$ as a functor to a space. The pointwise formula for its left Kan extension sends $Y$ at $b$ to the right side of \eqref{eq:boundary-fibre}. The equivalence between spaces over $G_D$ and space-valued local systems on $G_D$ identifies this left Kan extension with the space $\hocolim Y$ equipped with its projection to $G_D$. Its value at $b$ is exactly the homotopy fibre. This proves the formula without a false identification with one diagram value.

For the categorical assertion, represent an $\infty$-category by its complete Segal space and represent $G_D$ by a constant simplicial space. Pullback over $G_D$ is then levelwise. Groupoid completion is realization, and colimits of spaces commute with base change over a space. Thus groupoid completion commutes with this particular pullback, giving the displayed equivalence. No assertion that groupoid completion preserves arbitrary categorical pullbacks is being made.

Two examples expose the information retained by the formula. For the three coordinate lines in $\PP^2$, the closed divisor has homotopy type
\[
 D\simeq S^1\vee S^2\vee S^2\vee S^2,\qquad \Delta_D\simeq S^1.
\]
Indeed it is a graph of three sphere spaces joined by three point intersections; collapsing a spanning tree and arcs joining the attachment points gives this wedge. The fibre of $q_D$ is its infinite cyclic cover, a tree of spheres and hence a countable wedge of $S^2$'s. It is neither one line nor one point intersection. Next replace any pair by its product with $BH$ for a discrete group $H$. All component sets $K_q$ are unchanged, whereas every diagram value, total homotopy type, and fibre acquires the factor $BH$. In particular a groupoid Grothendieck construction retaining automorphisms computes a homotopy-orbit object, not the coarse dual complex. This also explains why taking components and retaining component groupoids are different operations.
\end{proof}

\begin{proposition}[What a boundary maquette can model]
\label{P4b}
\status{Proved}. If every connected component of each closed normalized branch intersection has contractible stack homotopy type, then $q_D$ is a homotopy equivalence. In contrast, if $X$ is a smooth compact complex manifold of pure dimension $n>0$, or a proper smooth analytic orbifold with finite stabilizers, and $D\ne\varnothing$, then
\[
 U_\infty\not\simeq\Delta_D.
\]
This obstruction includes ineffective finite generic stabilizers.
\end{proposition}
\begin{proof}
Each flag component is a component of a finite covering of a normalized deepest intersection. A covering of a contractible homotopy type has contractible connected components. Thus every $Y_D(a)$ is contractible, and the objectwise equivalence $Y_D\to *$ induces
\[
 D\simeq\hocolim Y_D\simeq\hocolim *\simeq\Delta_D.
\]
The hypothesis is about the closed normalized intersections, not merely their open strata or their coarse moduli spaces.

For the second assertion, rounding corners turns the nonempty boundary of $X^{\log}$ into a closed oriented real $(2n-1)$-manifold. Its orientation is induced by the complex orientation and the outward normal. At least one component therefore has a nonzero rational fundamental class, so
\[
 H^{2n-1}(U_\infty;\Q)\ne0.
\]
For a proper finite-isotropy orbifold the same conclusion follows from its rational orientation class. Locally the finite groups preserve the complex orientation, and hence the induced boundary orientation. The map from its stack homotopy type to its coarse orbifold induces an isomorphism in rational cohomology, since the higher rational cohomology of each finite stabilizer vanishes. Thus the compact oriented coarse orbifold's top class remains nonzero on the stack. Ineffective finite inertia does not remove it.

On the other hand, a dual simplex has at most $n$ vertices, so $\dim\Delta_D\le n-1$. Its cohomology in degree $2n-1$ vanishes. This proves the non-equivalence. The empty-boundary case is excluded, and the two assertions concern different maps: contractibility of closed intersections controls $D$, not the angular boundary $U_\infty$.

On MPL~149, fol.~66, scan~133, the final $U_\infty$ is legible and is retained in the reading register. The corresponding proposed maquette equivalence is therefore refuted under the compact hypotheses just proved. Replacing its target by $D$ under the stated contractibility hypothesis is a mathematical correction, not an emendation silently substituted into the transcription.
\end{proof}

\begin{theorem}[Admissible modifications and the scope of independence]
\label{P4c1}
\status{Proved}. Let $Z\subset D$ be a closed smooth centre having normal crossings with $D$, and put
\[
 X'=\operatorname{Bl}_Z X,\qquad D'=(f^{-1}D)_{\mathrm{red}}.
\]
Then $\Delta_{D'}\simeq\Delta_D$. Each connected component of the centre is treated separately: a stratum closure gives a stellar subdivision; a properly contained admissible centre gives a complex strongly deforming onto the old one. The local statement applies to the analytic DM scope above.

Consequently the coarse boundary homotopy type is independent of a smooth proper algebraic normal-crossings compactification of a fixed algebraic open variety, and likewise of a smooth proper algebraic DM stack in characteristic zero. For compact complex manifolds, the same conclusion holds when the specified identification of the open complements extends to a bimeromorphic map of their compactifications. The latter extension is a genuine hypothesis.
\end{theorem}
\begin{proof}
Work one connected component of $Z$ at a time; distinct components are disjoint and their blow-ups commute. Admissibility means that at every point there are simultaneous coordinates
\[
 D=\{x_1\cdots x_k=0\},\qquad Z=\{x_i=0:i\in I\}.
\]
Put $c=|I|$, $r=|I\cap\{1,\ldots,k\}|$, and $s=c-r$. If $s>0$, every local branch intersection leaves at least one nonboundary coordinate imposed by $Z$ free, so no branch-intersection component is contained in $Z$. If $s=0$ at some point, $Z$ agrees there with a codimension-$c$ branch intersection. Closedness and irreducibility of the connected smooth centre imply by analytic continuation that it is that stratum closure globally. Therefore a connected nonstratal admissible centre has $s>0$ everywhere and contains no deeper stratum component. In cone-extension notation this is $\Delta_Z^0=\varnothing$. The assertion would not follow from smoothness of the centre without the simultaneous normal-crossings coordinates.

Let $R$ be the set of the $r$ branches containing $Z$. Over $Z$, the exceptional divisor is $\PP(N_{Z/X})$, with fibre $\PP^{c-1}$. Intersecting its fibre with the strict transforms of $j$ distinct containing branches cuts out the corresponding $j$ coordinate hyperplanes. If $s=0$, the intersection with all $r=c$ containing branches is empty, whereas all proper subsets give the expected nonempty projective coordinate subspaces. Thus the star of the face corresponding to $Z$ is replaced by the cone from the exceptional vertex over its boundary, together with its unchanged link. This is stellar subdivision. It is equivariant under permutations of the containing branches and therefore descends to the generalized coarse dual complex.

If $s>0$, no old intersection component is contained in $Z$, so every old face remains. The new faces carry an exceptional vertex $v$ and are indexed by components of intersections with $Z$. For an old incidence $\sigma$ meeting $Z$,
\[
 D_\sigma\cap Z=D_{\sigma\cup R}\cap Z
\]
with identical connected-component labels. Over each such component, intersecting the exceptional divisor with all containing strict transforms gives a $\PP^{s-1}$-bundle, which is nonempty. Hence the coface $v*(\sigma\cup R)$ exists. On a new simplex define
\[
 H_t\left(\lambda_vv+\sum_{i\in\sigma}\lambda_i i\right)
 =(1-t)\lambda_vv+\sum_{i\in\sigma}\lambda_i i
       +t\lambda_v b_R,
\]
where $b_R$ is the barycentre of $R$. Its support lies in the indicated coface. It fixes every old simplex and at $t=1$ lands in the old complex. The equality of component labels makes these formulas compatible with all face identifications, including when $D_\sigma\cap Z$ is disconnected. The barycentre is invariant under branch monodromy, so the contraction descends. This proves the required strong deformation retraction. Applying the construction to an entire étale Čech nerve proves the local stack statement, with no finite-inertia restriction.

The simultaneous-coordinate hypothesis and the three combinatorial cases are precisely those in Stepanov's lemma as stated in \cite[p.~307]{Payne2013}; its Remark~5.1 gives the deformation retraction. For varieties, weak factorization supplies a sequence of these admissible blow-ups and blow-downs; Payne's Theorem~1.1 gives the stronger simple-homotopy independence. For proper algebraic DM stacks use \cite[Theorems~1.1--1.2]{Harper2017}: its factorization also permits boundary root constructions, which change stabilizers but not coarse branch-incidence components. For compact complex manifolds, \cite[Theorem~0.3.1, especially its normal-crossings clause]{AKMW2002} supplies the required factorization of a bimeromorphic map. This citation is not being asserted as an analytic-stack factorization theorem.

Thuillier gives a different, non-Archimedean construction and independence argument. The verified references are Theorem~3.25, Proposition~4.5, and Theorem~4.6 of the 28 July 2006 preprint \cite{Thuillier2006}. The journal version is \cite{Thuillier2007}; its theorem numbering has not been conflated with that checked preprint. Nothing in either the algebraic or bimeromorphic result implies invariance under an arbitrary biholomorphism of analytic open complements, as the next example proves.
\end{proof}

\begin{example}[Two good analytic compactifications with different maquettes]
\label{P4c3}
\status{Refuted}. Analytic compactification-independence of the coarse dual complex fails even for good compactifications. The same analytic open surface $(\C^*)^2$ has good compactifications with dual complexes $S^1$ and a point, while both links at infinity are $T^3$.
\end{example}
\begin{proof}
The standard toric pair
\[
 X_T=(\PP^1)^2,\qquad D_T=X_T\setminus(\C^*)^2
\]
has four boundary curves arranged in a square. Thus $\Delta_{D_T}\cong S^1$. Its open strata are $(\C^*)^2$, copies of $\C^*$, and points, all aspherical. The local meridians inject into $\pi_1((\C^*)^2)=\Z^2$: at a vertex they are a signed basis, and along an edge a primitive generator. Hence the pair is good. Its oriented real blow-up is $T^2\times[0,1]^2$, whose rounded boundary is $T^3$.

Choose $0<\alpha<1$, put
\[
 g(z,w)=(\alpha z+w,\alpha w),
 \qquad H=(\C^2\setminus\{0\})/\langle g\rangle,
 \qquad E=\{w=0\}/\langle g\rangle.
\]
The linear map $g$ is a contraction for an equivalent norm: its spectral radius is $\alpha<1$. Its nonzero integer powers have no fixed points away from the origin, and contraction shows that the action is properly discontinuous. A closed fundamental annulus gives compactness of the quotient. Thus $H$ is a smooth compact Hopf surface, and
\[
 E\cong\C^*/\alpha^{\Z}
\]
is a smooth elliptic curve. This is the Hopf case in \cite[Theorem~2, p.~81; \S4(d), p.~88]{Ueda1979}.

On $w\ne0$ introduce the correctly scaled coordinates
\[
 u=\alpha z/w,\qquad v=w.
\]
Then $g(u,v)=(u+1,\alpha v)$. Writing $\ell=\log\alpha$, the functions
\[
 A=\exp(2\pi i u),\qquad B=v\exp(-\ell u)
\]
are invariant and define $H\setminus E\to(\C^*)^2$. Conversely choose a local logarithm of $A$, put $u=(2\pi i)^{-1}\log A$, $v=B\exp(\ell u)$ and $z=uv/\alpha$. Changing the logarithm changes $u$ by an integer and applies the corresponding power of $g$. This proves that the map is a biholomorphism.

It remains to verify goodness, rather than merely identifying the complement. Both $H\setminus E\cong(\C^*)^2$ and $E$ are aspherical. A small normal meridian can be represented upstairs by holding $z=z_0\ne0$ fixed and letting $w=\varepsilon\exp(2\pi it)$. The function $u$ is single-valued along this closed path, so $A=\exp(2\pi iu)$ has winding number zero. The factor $\exp(-\ell u)$ also has winding zero, while $v$ has winding one. Thus the meridian has winding vector $(0,1)$ in $(A,B)$, and its map $\Z\to\Z^2$ is injective. These are all open strata and all local meridian conditions, so the Hopf pair is good.

Its dual complex has one vertex. To identify its link, use the normal coordinate $t=w/z$ near the inverse image of $E$. Under $g$ it transforms as
\[
 t\longmapsto\frac{\alpha t}{\alpha+t},
\]
whose derivative at $t=0$ is $1$. Thus the normal line bundle descends with trivial transition on its fibres and is holomorphically trivial. Its unit circle bundle is $E\times S^1\cong T^3$, which is the boundary link. We have therefore obtained
\[
 \Delta_{D_T}\simeq S^1,\qquad\Delta_E\simeq *,
 \qquad L_T\cong L_H\cong T^3.
\]
The prescribed biholomorphism of open surfaces consequently cannot extend to a bimeromorphic map of these compactifications. The program \texttt{boundary\_analytic\_checks.js}, executed at $\alpha=1/2$, checks the affine coordinate transformation, the normal derivative, the primitive meridian vector, the two finite dual graphs, and the homology ranks of $T^3$. Its checks supplement the quotient and winding proofs above; they do not replace them.
\end{proof}

\begin{theorem}[Weight zero of the link and compact-support/top-weight comparison]
\label{P4c2}
\status{Proved}. Assume now that $X$ is connected, smooth and proper algebraic over $\C$, of dimension $n>0$, with $D\ne\varnothing$ a normal-crossings divisor. The same statements hold for a connected smooth proper algebraic DM stack with finite inertia. The algebraic pair equips the link cohomology with its mixed Hodge structure, and for every $k\ge0$,
\[
 H^k(\Delta_D;\Q)
 \xrightarrow{\ \sim\ } W_0H^k(D;\Q)
 \xrightarrow{\ \sim\ } W_0H^k(L;\Q).
\]
The maps are induced by $q_D$ and restriction to the link, and their composite is $q_L^*$. For $k\ge1$ there are natural isomorphisms
\[
 W_0H_c^k(U;\Q)\cong\widetilde H^{k-1}(\Delta_D;\Q),
\]
\[
 \operatorname{Gr}_{2n}^{W}H^{2n-k}(U;\Q)
 \cong\widetilde H_{k-1}(\Delta_D;\Q)\otimes\Q(-n).
\]
The reduced groups in degree zero are necessary. For disconnected $X$ the assertions are applied componentwise; in degree one the reduction is the cokernel of $H^0(X)\to H^0(D)$.
\end{theorem}
\begin{proof}
We first specify the Hodge input and its stack scope. For varieties the mixed Hodge structure on link cohomology is constructed through the local object $i^*Rj_*\Q_U$, where $i:D\hookrightarrow X$ and $j:U\hookrightarrow X$; see \cite[Definitions~2.7--2.8 and Propositions~2.9, 3.5]{DurfeeSaito1990}. The closed angular-face model identifies its global Betti cohomology with that of $L$. For algebraic stacks we use the mixed-Hodge-module formalism in \cite{Tubach2025}: Theorem~4.1.1 provides the six operations, purity and localization; Propositions~4.2.5--4.2.6 provide the proper DM and Betti comparisons; Propositions~4.3.5--4.3.6 and Corollary~4.3.7 give the weight compatibilities. In particular smooth proper DM cohomology is pure of its cohomological degree, proper cohomology has weights at most that degree, and localization and duality are compatible with the weight filtration. Finite inertia is essential to the proper rational-cohomology scope used here.

Use the proper normalized-flag resolution from P4a2. Every flag component is smooth and proper; its cohomology in degree $q$ is pure of weight $q$. The resulting cohomological spectral sequence has
\[
 E_1^{p,q}=\bigoplus_{a\in K_p}H^q(Y_D(a);\Q)
 \Longrightarrow H^{p+q}(D;\Q).
\]
Its weight-zero row is the ordinary component cochain complex of $K$. The differentials $d_r$ for $r\ge2$ between different rows vanish by weights: their source and target have distinct pure weights. Therefore the weight-zero part of the abutment is the cohomology of that row, namely $H^k(\Delta_D;\Q)$. The identification is induced by $q_D$, so it also respects products.

Localization gives a triangle
\[
 i^!\Q_X\longrightarrow\Q_D\longrightarrow i^*Rj_*\Q_U
 \longrightarrow i^!\Q_X[1]
\]
and hence an exact sequence of mixed Hodge structures
\[
 H_D^k(X)\longrightarrow H^k(D)\longrightarrow H^k(L)
 \longrightarrow H_D^{k+1}(X).
\]
Proper duality identifies
\[
 H_D^a(X;\Q)\cong H^{2n-a}(D;\Q)^\vee(-n).
\]
Since the cohomology group on the right before duality has weights at most $2n-a$, the support group has weights at least $a$. Thus its weight-zero part vanishes for $a>0$. The endpoint $a=0$ requires a separate argument: $H_D^0(X)=0$, since a locally constant function supported on $D$ is zero, as $D$ contains no component of $X$. Therefore both adjacent support terms have zero weight-zero part for every $k\ge0$, including $k=0$. Strictness of morphisms of mixed Hodge structures proves $W_0H^k(D)\cong W_0H^k(L)$. Compatibility of the diagrams in P4a2 identifies its composite with $q_L^*$.

The compact-support sequence is
\[
 H^{k-1}(X)\longrightarrow H^{k-1}(D)
 \longrightarrow H_c^k(U)\longrightarrow H^k(X).
\]
For $k\ge2$, purity of the two groups on $X$ makes their weight-zero parts vanish, giving the first compact-support formula. For $k=1$ the exact weight-zero sequence instead gives
\[
 W_0H_c^1(U)=\operatorname{coker}\bigl(H^0(X)\to H^0(D)\bigr)
 =\widetilde H^0(\Delta_D),
\]
because $X$ is connected and $D$ is nonempty. Also $H_c^0(U)=0$ by the same injectivity of $H^0(X)\to H^0(D)$. Finally Poincaré duality identifies
\[
 H^{2n-k}(U)\cong H_c^k(U)^\vee(-n).
\]
Dualizing the weight-zero formula gives the displayed top-weight homology formula, with the Tate twist recording its Hodge weight. For varieties these comparisons agree with \cite[Proposition~4.2 and Theorem~4.4]{Payne2013}. No such argument assigns these algebraic weight comparisons to arbitrary analytic compactifications; P4c3 explicitly prevents that extension.
\end{proof}

\begin{proposition}[Stable graphs, tropical links, and automorphisms]
\label{P4d1}
\status{Proved}. Assume $2g-2+n>0$. The coarse dual complex of the boundary of $\overline{\mathcal M}_{g,n}$ is the length-one tropical moduli space $\Delta_{g,n}$. It is obtained from stable graphs with at least one edge by ordinary simplex quotients and zero-length edge contractions. It is not, in general, the classifying space of the graph contraction category retaining graph automorphisms, nor is it the full finite-length tropical cone.
\end{proposition}
\begin{proof}
A stable graph $\Gamma$ is connected, has labelled legs $1,\ldots,n$, a genus $g(v)\ge0$ at every vertex, and satisfies
\[
 g=\sum_vg(v)+b_1(\Gamma),\qquad
 2g(v)-2+\operatorname{val}(v)>0.
\]
Its boundary stratum is
\[
 \left[\prod_{v\in V(\Gamma)}
 \mathcal M_{g(v),\operatorname{val}(v)}
 \big/\operatorname{Aut}(\Gamma)\right],
\]
with the half-edges regarded as marked points before gluing. These factors are connected, and their proper counterparts give the closed normalized incidence spaces. Each edge corresponds to a node, with a local smoothing coordinate. The normal-crossings branches are precisely the vanishing node parameters, and their monodromy is the permutation of edges induced by graph isomorphisms. Curve automorphisms fixing all nodes may still contribute inertia, but they do not produce additional coarse branches.

For a graph with $e>0$ edges attach
\[
 \sigma_\Gamma=
 \{(\ell_a)_{a\in E(\Gamma)}\in\R_{\ge0}^{e}:
                   \sum_a\ell_a=1\}.
\]
Take its ordinary topological quotient by $\operatorname{Aut}(\Gamma)$. If a collection of lengths becomes zero, contract those edges; a contracted loop increases the genus of its vertex, and more generally a contracted connected subgraph contributes its first Betti number to the merged vertex genus. These are exactly the incidence face maps of node smoothing. Consequently the glued simplex quotient is the ordinary symmetric-delta realization of the boundary. Its full-flag subdivision is the construction of P4a2. This proves the identification with $\Delta_{g,n}$; compare \cite[Definition~5.2, Lemma~5.4, Corollaries~5.6--5.7 and Theorem~5.8]{CGP2021}, and the cone-complex and tropical identifications in \cite[Theorem~1.2.1, Corollary~7.2.2 and \S7.3]{ACP2015}.

The distinction from homotopy quotients is visible in the smallest stable example. For $(g,n)=(1,1)$ the only positive-edge graph has one genus-zero vertex, one labelled leg and one loop. Its automorphism group is $C_2$, interchanging the loop's half-edges. This group fixes the unique edge and hence acts trivially on its one-point length-one simplex. Thus $\Delta_{1,1}$ is a point. The positive-edge graph category retaining automorphisms instead has one object with automorphism group $C_2$, so its classifying space is $BC_2$. Its first integral homology is $\Z/2$, and it is not a point. The normalized bar differential in degree two is multiplication by $2$, as also checked by \texttt{boundary\_checks.py}. This is an explicit failure of substituting a groupoid Grothendieck construction for the coarse realization.

Allowing all lengths with no normalization gives the finite-length tropical cone: every nonzero point is a positive scalar times its unique length-one representative. The no-edge graph is the cone vertex, reached by shrinking all edges. It belongs to this cone, not to its link. For $(0,3)$ the compactified moduli space has empty boundary and $\Delta_{0,3}$ is empty. The types $(0,0),(0,1),(0,2),(1,0)$ do not satisfy the stability inequality and are not exceptions to be filled by an artificial cone vertex.
\end{proof}

\begin{theorem}[Genus-zero homotopy type and exact finite checks]
\label{P4d2}
\status{Known}. For $n\ge4$,
\[
 \Delta_{0,n}\simeq\bigvee^{(n-2)!}S^{n-4}.
\]
In particular the face vectors for $n=4,5,6$ are respectively
\[
 (3),\qquad(10,15),\qquad(25,105,105),
\]
and their reduced integral homology is respectively $\Z^2$ in degree zero, $\Z^6$ in degree one, and $\Z^{24}$ in degree two, with all other reduced groups zero. These finite coarse-complex computations do not verify the infinite posets $P_k$ considered in Part~5.
\end{theorem}
\begin{proof}
In genus zero every stable graph is a tree and every vertex has genus zero. Root the tree at the labelled leaf $n$. Every internal edge determines the set $A\subset\{1,\ldots,n-1\}$ of leaves on its side away from the root. Stability gives
\[
 2\le|A|\le n-2.
\]
Two such sets are either nested or disjoint. Thus a simplex determines a laminar family of these nontrivial proper subsets.

Conversely, let $\mathcal F$ be such a laminar family. Adjoin the full set $\{1,\ldots,n-1\}$ and every singleton. Connect each member other than the full set to its smallest strict superset in this enlarged family; laminarity makes that parent unique. Regard the singletons as labelled leaves and attach leaf $n$ to the full-set vertex. A nonsingleton vertex has at least two children: one child cannot contain all its labels without being the vertex itself. With its parent edge, or with the added root leaf for the full-set vertex, it therefore has valence at least three. This constructs a stable labelled tree. Its internal edges are exactly the proper nonsingleton members of $\mathcal F$, and deleting one member contracts the corresponding edge. The two constructions are inverse on simplices and their face maps. Also a genus-zero stable tree has no nontrivial automorphism fixing every labelled leaf: its edge splits determine every internal edge and then every vertex. Hence there is no remaining graph quotient in this case.

The resulting complex is exactly the Whitehouse complex in the unnumbered definition of \S2, preprint p.~3, of \cite{Readdy2005}, after relabelling the root leaf. Readdy's Theorem~4.7 gives its homotopy equivalence to $(n-2)!$ spheres of dimension $n-4$; Theorem~5.2 gives the associated enumerative information. This identifies every simplex, face, indexing convention and dimension needed to apply that theorem. The comparison with the rooted tree/nested-set model is also made through stellar subdivisions in \cite[Theorem~3.1 and Corollary~3.2]{Feichtner2006}. In particular the assertion is about homotopy type, not merely a numerical Euler-characteristic calculation. For $n=4$, the convention that a wedge of two $S^0$'s is three points gives the correct boundary.

Here are independently executed integral checks of the small cases. The program \texttt{boundary\_checks.py} first lists all subsets of $\{1,\ldots,n-1\}$ of sizes $2,\ldots,n-2$ and all pairwise compatible families; laminarity makes pairwise compatibility sufficient. It sorts the vertices to orient each simplex and forms the integer simplicial boundary matrices with alternating signs. For $n=4$ there are three vertices and no edges. For $n=5$ the graph has ten vertices and fifteen edges and is connected, so its first Betti number is $15-10+1=6$.

For $n=6$ the program obtains twenty-five vertices, 105 edges and 105 triangles, with no higher simplices, and verifies $\partial_1\partial_2=0$. A spanning tree has twenty-four edges. Projection of integral cycles onto the remaining eighty-one edge coordinates is an isomorphism from the integral cycle group to $\Z^{81}$: the omitted tree coefficients are uniquely recovered from the zero-boundary condition. Under this isomorphism $\partial_2$ is an $81\times105$ integer matrix. Its Smith normal form has eighty-one nonzero diagonal entries, all equal to $1$. Consequently the triangle boundaries generate the full cycle lattice, so $H_1=0$ integrally, not merely rationally. The kernel of $\partial_2$ is free of rank $105-81=24$; because there are no three-simplices, it is $H_2$. These calculations prove the stated integral groups. They do not by themselves prove simple connectivity; the homotopy assertion uses the independently identified Whitehouse-complex theorem.
\end{proof}

\begin{proposition}[What the coarse maquette forgets]
\label{P4e}
\status{Proved}. The coarse dual complex does not determine the homotopy type of $U_\infty$, even among good smooth projective SNC pairs. It forgets both the topology of the normalized strata and the angular normal-bundle data. Conversely, homeomorphic links of analytic compactifications need not determine the same coarse dual homotopy type.
\end{proposition}
\begin{proof}
For the first assertion within the good projective category, fix an elliptic curve $B$ and a point $p\in B$. For a smooth projective curve $C_g$ of genus $g\ge1$, take
\[
 X_g=C_g\times B,\qquad D_g=C_g\times\{p\},
 \qquad U_g=C_g\times(B\setminus\{p\}).
\]
The open strata are $U_g$ and $D_g$. Both are aspherical: a positive-genus closed curve is aspherical, and a punctured elliptic curve retracts onto a wedge of two circles. The unique local meridian maps to the puncture loop in
$\pi_1(B\setminus\{p\})=F_2$. This loop is the commutator of a free basis, hence nontrivial and of infinite order. Its cyclic subgroup injects, so the pair is good. Its boundary has one smooth component and therefore $\Delta_{D_g}$ is a point. The normal line bundle is the fixed tangent line $T_pB$ over $C_g$, so
\[
 L_g\cong C_g\times S^1.
\]
For $g=1$ and $g=2$ these links have first Betti numbers $3$ and $5$, respectively, and cannot be homotopy equivalent. Thus even a one-point maquette loses essential information in the good projective setting.

Normal-bundle information can be isolated while fixing the boundary's own homotopy type. In $\PP^2$, take first a line and then a smooth conic. Both boundary curves are isomorphic to $\PP^1$, and both coarse dual complexes are points. For a degree-$d$ plane divisor, the normal line bundle has degree its self-intersection $d^2$; therefore the two normal bundles here are $\mathcal O_{\PP^1}(1)$ and $\mathcal O_{\PP^1}(4)$. The link is their unit circle bundle. In the homotopy exact sequence
\[
 \pi_2(S^2)=\Z\longrightarrow\pi_1(S^1)=\Z
 \longrightarrow\pi_1(L)\longrightarrow0,
\]
the first map is multiplication by the Euler number. The two link fundamental groups are consequently trivial and $\Z/4$. More precisely the first link is $S^3$ and the second a lens space $L(4,1)$, up to the convention for orientation. The elementary integer cokernel calculation is included in \texttt{boundary\_checks.py}. These plane-curve pairs are being used to isolate normal data; no goodness claim about them is needed.

The diagrams in P4a2 retain exactly the kinds of information discarded in these examples: the homotopy types of the normalized flag components, their angular torus bundles and characteristic classes, the gluing and monodromy maps, the higher coherence of those maps, and stack inertia. Replacing every value by its component label gives the coarse realization and forgets this data. Hence there is no reconstruction of the full link from that realization alone, even when the crossing category happens to be an ordinary category.

For the converse, P4c3 gives homeomorphic links $T^3$ of the same analytic open surface with dual complexes $S^1$ and a point. Thus neither the link homotopy type nor an unrestricted analytic identification of the open complement recovers the coarse maquette. The algebraic and bimeromorphic independence statements of P4c1 remain valid in their stated scopes.
\end{proof}

\section{Codimension levels and modular dimension levels}
\label{sec:levels}

The manuscript uses both a filtration by normal-crossing codimension and
a filtration by modular dimension. They have different index posets and
connectivity estimates. A map is \(r\)-connected here when its homotopy
fibres are \((r-1)\)-connected; it induces isomorphisms on \(\pi_i\) for
\(i<r\) and an epimorphism on \(\pi_r\). A \(k\)-type equivalence is an
equivalence after \(\tau_{\leq k}\).

\subsection{The diagram (246), its coherence, and the sharp cutoff}

Let \(L=\partial X^{\log}\) be the boundary of the real oriented blow-up
of an NC pair. Define \(L_{\leq r}\) as the inverse image of the locus
where at most \(r\) local branches meet. Thus the truncation is an actual
subspace, with every incidence flag retained. For a proper pair a rounded
collar identifies \(L\) with a deleted neighbourhood of the divisor in
the open part. For a nonproper pair we use the resolved boundary itself;
an arbitrary end of the open space need not be such a collar.

\begin{proposition}[The two-level reading of (246)]\label{P5a1}
\status{Refuted.}
The bicategorical colimit in groupoids of the depth-one/depth-two
incidence span displayed in (246) need not be \(\Pi_1(L)\).
For the three coordinate hyperplanes in \(\C^3\), that colimit has vertex
group \(\Z^4\), whereas \(\pi_1(L)=\Z^3\). The depth-three object and its
flag triangles remove the extra generator.
\end{proposition}

\begin{proof}
Polar coordinates give
\[
 L=(S^1)^3\times
 \{(r_1,r_2,r_3)\in\R_{\geq0}^3:\min r_i=0\}.
\]
The radial factor contracts to the origin, so \(L\simeq(S^1)^3\).
Depth at most two deletes that origin. Dividing by the positive radial
coordinate identifies the remaining radial factor with the boundary
of a two-simplex. Hence \(L_{\leq2}\simeq(S^1)^3\times S^1\).

This also audits the actual two-level incidence diagram. Its indices
are the three singleton and three two-element subsets of
\(\{1,2,3\}\), with six inclusions, forming a hexagon. Each total-link
space is \((S^1)^3\) up to homotopy and every map induces the identity
on its three coordinate meridians. The bicategorical groupoid colimit
is the fundamental groupoid of the homotopy colimit of these classifying
spaces. The hexagon therefore supplies an additional infinite cyclic
generator. A spanning-tree presentation has one index generator commuting
with the three meridians.

Adjoining \(\{1,2,3\}\) and all composable flags gives the barycentric
subdivision of a filled two-simplex. The six triangular relations
eliminate the six non-tree edges of its enlarged graph.
The program \texttt{levels\_colimit.py} computes both presentations by
single-occurrence Tietze eliminations. Its outputs are
\((V,E,F)=(6,6,0)\) and \((7,12,6)\), with one and zero surviving
index generators respectively.

The strict colimit of the corresponding one-object groups would already
be \(\Z^3\). That strict colimit discards the path and natural-isomorphism
data required for descent and does not interpret the bicategorical
colimit in (246). The fully coherent flag diagram includes the
depth-three triangles and gives the correct answer.
\end{proof}

The marginal reference to \(2\)-fl\`eches on folio 76, scan 152,
supports recording the natural isomorphisms on triple flags. It does
not make an unwritten triangular relation automatic. The neighbouring
purity sentence cannot justify deleting depth three: within the
resolved boundary, depth three has real codimension two.

\begin{theorem}[Sharp codimension truncation]\label{P5a2}
\status{Proved.}
For every \(r\geq1\), the inclusion
\[
                         L_{\leq r}\longrightarrow L
\]
is \((r-1)\)-connected. Consequently depths at most \(k+2\), with the
full flag diagram, determine \(\tau_{\leq k}L\). The uniform bound is
optimal: for positive integral cutoffs \(r\), depth at most \(r\)
determines the \(k\)-type for every NC pair if and only if \(r\geq k+2\).

The assertion holds for analytic DM stacks on classifying spaces.
If every stratum is aspherical, the total normal-torus links in this
derived reconstruction can be replaced by classifying spaces of
their fundamental groupoids. This replacement does not require
meridian injectivity into the open-stratum group.
\end{theorem}

\begin{proof}
At a point of depth \(a\), a chart of \(X^{\log}\) has \(a\)
nonnegative radial coordinates and \(a\) circle coordinates. Within
its boundary the locus where those radial coordinates all vanish
has real codimension \(a-1\). The removed strata, of depths at least
\(r+1\), therefore have real codimension at least \(r\) in \(L\).

Triangulate the corner charts compatibly, or round their corners and
use the resulting locally flat strata. Relative general position moves
a map of a compact polyhedron of dimension less than \(r\) off the
removed union while fixing an already disjoint closed subpolyhedron.
Applied to spheres and nullhomotopies, this gives isomorphisms on
\(\pi_i\) for \(i<r-1\) and surjectivity on \(\pi_{r-1}\).
For noncompact manifolds it is applied to the compact image of each
map; no stable-end assertion is needed. For example, for
\(X=\C^2,D=\{z_1=0\}\), the normal boundary is
\(\C\times S^1\simeq S^1\), while the full end of
\(X^*=\C^*\times\C\cong S^1\times\R^3\) has cross-section
\(S^1\times S^2\). Thus even these homotopy types can differ.

For a stack choose an \'{e}tale atlas and its entire \v{C}ech nerve.
The manifold result applies to each nerve term. Geometric realization
preserves a fixed connectivity bound for maps of spaces. One proof
uses relative CW models: their relative cells have the required
minimum dimension, and realization adds a nonnegative simplicial
dimension. Equivalently, the appropriate left class of the
connected/truncated factorization system is closed under homotopy
colimits. Thus the termwise inclusions give the same bound on
classifying spaces. This argument includes nonproper analytic DM
stacks with infinite discrete isotropy, rather than assuming finite
orbifold charts.

For necessity take \(m=r+1\) coordinate hyperplanes in \(\C^m\).
The radial calculation above gives
\[
 L\simeq(S^1)^m,\qquad
 L_{\leq r}\simeq(S^1)^m\times S^{m-2}
                    =(S^1)^{r+1}\times S^{r-1}.
\]
For \(r=1\) this changes components; for \(r\geq2\) it changes
\(\pi_{r-1}\). Whenever \(r\leq k+1\), the \(k\)-types differ.
All these sharpness examples already have aspherical strata and
injective full meridian lattices. The supplied exact boundary-matrix
calculations verify the simplex boundary homology for dimensions
\(0,\ldots,4\); the radial proof establishes the assertion in every dimension.

Finally, a normal-torus bundle over an aspherical stratum is
aspherical by its homotopy exact sequence, as are the iterated flag
links. Objectwise replacement by \(B\Pi_1\) is consequently an
equivalence of coherent diagrams, and remains so under homotopy
colimits. This concerns a derived diagram of spaces; replacement of
the exit category by its ordinary homotopy category is a different
operation.
\end{proof}

\begin{corollary}[The corrected fundamental-group conclusion]\label{P5a3}
\status{Proved.}
Coherent fundamental-groupoid data from depths one, two, and three
always determine \(\Pi_1(L)\). Depths one and two suffice if there
are no triple crossings. If the collar inclusion
\(L\to X^*=X\setminus D\) is a \(1\)-type equivalence, the conclusion
is \(\Pi_1(X^*)\), and need not be \(\Pi_1(X)\).
\end{corollary}

\begin{proof}
Apply Theorem~\ref{P5a2} with \(k=1\). The fundamental-groupoid
functor is the left adjoint from spaces to one-types and therefore
preserves colimits, interpreted as bicategorical colimits in groupoids.
Consequently the full coherent diagram of fundamental groupoids always
computes \(\Pi_1(L)\), without an asphericity assumption on its pieces.
Without triple crossings, \(L_{\leq2}=L\). Compose with the assumed
equivalence to \(X^*\).
For the last distinction take \((X,D)=(\C,\{0\})\): its collar is
\(S^1\simeq\C^*\), but \(X\) is simply connected. Thus the unstarred
\(X\) visibly written in (247), scan 152, needs a mathematical
correction under this hypothesis; it is not silently transcribed
as \(X^*\).
\end{proof}

\subsection{Dimension filtration and its actual homotopy fibre}

Let \((g,n)\) be stable, \(d=3g-3+n\), and let \(\Gamma_{g,n}\)
fix the \(n\) labels individually. A multicurve means an isotopy
class of a finite family of pairwise disjoint, essential,
nonperipheral, mutually nonisotopic simple closed curves.
The empty multicurve is explicitly allowed. The stratum associated
to \(\sigma\) has dimension \(d-|\sigma|\). For every \(k\geq0\), put
\[
 P_k=\{\sigma:|\sigma|\geq d-k\},\qquad
 \sigma\leq\tau\ \Longleftrightarrow\ \sigma\subseteq\tau .
\]
Let \(W_k\) be the union, inside the open moduli stack, of deleted
tubes around strata of dimension at most \(k\). A precise model
requires the curves of some such \(\sigma\) to have hyperbolic
length less than a fixed sufficiently small \(\epsilon\).
For \(k\geq d\), the open stratum itself is retained, and
\(W_k=\mathcal M_{g,n}\).

\begin{theorem}[Fibre of the modular dimension filtration]\label{P5b1}
\status{Proved.}
For every stable \((g,n)\) and every \(k\geq0\), the homotopy fibre of
\[
                          BW_k\longrightarrow B\mathcal M_{g,n}
\]
is homotopy equivalent to \(|P_k|\). All stack stabilizers are retained.
For \(k\geq d\) the poset has the minimum element \(\varnothing\),
so its realization is contractible.
\end{theorem}

\begin{proof}
Work in marked Teichm\"uller space \(\mathcal T_{g,n}\).
Choose \(\epsilon\) so that intersecting essential geodesics cannot
both have length less than \(\epsilon\), as supplied by the collar
inequality. Define
\[
 V_\sigma=\{x\in\mathcal T_{g,n}:
               \ell_\gamma(x)<\epsilon\text{ for all }\gamma\in\sigma\}.
\]
Extend \(\sigma\) to a pants decomposition. Fenchel--Nielsen
coordinates express \(V_\sigma\) as a product of \((0,\epsilon)\)
and \((0,\infty)\) length intervals and real twist lines, so it is
contractible; see \citet[Chapter 1, Theorem 1.1]{Harer1988}.
The convention \(V_\varnothing=\mathcal T_{g,n}\) is essential.

Set \(V_k=\bigcup_{\sigma\in P_k}V_\sigma\). A finite intersection
of these open sets is empty for an incompatible union of curves,
and otherwise is \(V_{\cup\sigma}\); the union again belongs to
\(P_k\). This numerable good cover is closed under nonempty finite
intersections. Its homotopy-colimit construction gives an
equivariant span of underlying weak equivalences
\[
 V_k\ \longleftarrow\
 \hocolim_{\sigma\in P_k^{\mathrm{op}}}V_\sigma
 \ \longrightarrow\ |P_k^{\mathrm{op}}|.
\]
For the left arrow, the augmented \v{C}ech indexing nerve at a
point has a largest object: the full, finite system of its short
curves. It is therefore contractible. For the right arrow use the
contractibility of each \(V_\sigma\). Opposite posets have the same
order complex up to its canonical reversal.

The same length and twist coordinates describe a deleted pinching
tube: length gives the radial coordinate of the neck, twist its
angular coordinate, and the remaining marked component coordinates
its base. Sufficiently small bounds admit contractible refinements.
The span thus computes the specified tubular model of \(W_k\);
compare the boundary construction of
\citet[Chapter 3, Lemma 3.6]{Harer1988}.

Finally use the Borel constructions
\[
 B\mathcal M_{g,n}\simeq
 E\Gamma_{g,n}\times_{\Gamma_{g,n}}\mathcal T_{g,n}
 \simeq B\Gamma_{g,n},\qquad
 BW_k\simeq E\Gamma_{g,n}\times_{\Gamma_{g,n}}V_k .
\]
The second space has homotopy fibre \(V_k\) over \(B\Gamma_{g,n}\).
Contractibility of Teichm\"uller space proves the assertion for
the displayed map to \(B\mathcal M_{g,n}\).
A coarse quotient would discard stabilizer classifying spaces,
including ineffective inertia, and is not the construction used.
For \(k\geq d\), \(V_\varnothing\) already makes \(V_k\) the whole
Teichm\"uller space and verifies the endpoint directly.
\end{proof}

This poset remembers actual isotopy classes on a marked surface.
It is infinite and is not the finite poset of stable labelled
trees in a coarse boundary complex. Thus the finite calculations
for \(\Delta_{0,5}\) and \(\Delta_{0,6}\) do not alone verify this
homotopy fibre.

\begin{theorem}[Genus-zero connectivity at every level]\label{P5b2}
\status{Proved.}
For \(n\geq3\) and \(k\geq0\), the realization \(|P_k|\) for the
\(n\)-punctured sphere is \((k-1)\)-connected.
Consequently \(BW_k\to\mathcal M_{0,n}\) is \(k\)-connected.
For \(k\geq n-3\), the fibre is contractible.
\end{theorem}

\begin{proof}
The endpoint was proved above and includes \(n=3\).
Suppose \(0\leq k<d=n-3\). Let \(K=\mathcal C_{0,n}\) be the
unmodified disjointness curve complex. It is pure of dimension
\(d-1\), and the genus-zero case of
\citet[Chapter 4, Theorem 4.1]{Harer1988} identifies its homotopy
type as a wedge of \((d-1)\)-spheres.

For a face \(\sigma\) with \(q\) vertices, cutting along its curves
gives punctured spheres with \(n_v\) special points, and
\[
              \operatorname{Lk}_{K}(\sigma)
                       =\mathop{*}_{v}\mathcal C_{0,n_v}.
\]
The \(n_v=3\) factors are empty complexes, the join units \(S^{-1}\).
All other factors have the sphere dimensions just stated.
Since \(\sum_v(n_v-3)=d-q\), this link is
\((d-q-2)\)-connected, with homotopy dimension \(d-q-1\).
This establishes the homotopy Cohen--Macaulay property of all links.

Here is a rank-selection proof valid for the infinite complex.
Use the weak CW realization of the barycentric subdivision of \(K\).
Successively delete vertices corresponding to faces of sizes
\(1,2,\ldots,d-k-1\). When a face of size \(q\) is deleted, its
remaining link consists of proper cofaces and is the barycentric
subdivision of \(\operatorname{Lk}_{K}(\sigma)\). It is at least
\((k-1)\)-connected because \(q\leq d-k-1\).
The pre-deletion complex is obtained from the post-deletion complex
by coning this link. That inclusion is \(k\)-connected and
preserves homotopy through degree \(k-1\).
Vertices of the same rank are incomparable; their open stars are
disjoint, and the cones may be attached simultaneously.
This also works for infinitely many vertices: relative CW
connectivity is preserved under the coproduct of these cone
attachments, or one can apply the argument to the finite
subcomplex traversed by a compact homotopy in the weak CW topology.
The initial complex is \((d-2)\)-connected and the final one is
\(|P_k|\). For \(k=0\), its required nonemptiness follows directly
from a pants decomposition.

This is the homotopical rank-selection mechanism of
\citet[Corollary 4.4]{BjornerWachsWelker2009}, with the infinite
case justified explicitly instead of inferred from a finite-poset
statement.
\end{proof}

Here \(\mathcal C_{0,4}\) is an infinite discrete set; using a
modified Farey complex changes the question. The empty multicurve
also cannot be omitted. For example, at \((n,k)=(6,3)\) omission
would leave \(\mathcal C_{0,6}\), instead of the contractible fibre.

\begin{proposition}[Independent checks at \(n=5,6\)]\label{P5b3}
\status{Proved.}
The fibres for \((n,k)=(5,1)\) and \((6,2)\) are respectively
connected and simply connected by an independent argument on the
actual configuration spaces. The fibre for \((6,3)\) is contractible.
\end{proposition}

\begin{proof}
Normalize three points to \(0,1,\infty\). For \(d=n-3\),
\[
 \mathcal M_{0,n}
 =\operatorname{Conf}_{d}(\C\setminus\{0,1\})
 =\C^d\setminus\bigcup\{z_i=0,\ z_i=1,\ z_i-z_j=0\}.
\]
Let \(\mathcal L\) be these affine linear forms, and put
\[
 f(z)=\sum_i |z_i|^2+\sum_{\ell\in\mathcal L}|\ell(z)|^{-2}.
\]
The function is proper: the first term controls infinity and the
others diverge along every deleted hyperplane. If \(a_\ell\)
denotes the coefficient vector of \(\ell\), its Levi matrix is
\[
                  I_d+\sum_{\ell\in\mathcal L}
                            |\ell(z)|^{-4}a_\ell a_\ell^*,
\]
which is positive definite. A generic real linear perturbation
leaves this matrix unchanged and makes \(f\) Morse.
The critical points are the inverse image of a regular value
of its real gradient. That gradient is semialgebraic, so its
zero-dimensional regular inverse image is finite.
The perturbed function remains proper and bounded below.

At a critical point the real Hessian satisfies
\(Q(v)+Q(Jv)>0\) for nonzero \(v\). A negative subspace of
dimension larger than \(d\) would intersect its \(J\)-translate
nontrivially, contradicting that inequality.
Every Morse index is therefore at most \(d\).
A regular sublevel above all critical values is a compact
\(2d\)-manifold with handle indices at most \(d\). Reversing
the handles gives a relative decomposition from its boundary
with indices at least \(d\). Its boundary inclusion is consequently
\((d-1)\)-connected. For the precise handle statements see
\citet[Theorems 3.2 and 3.5, Remark 3.3, Theorem 7.2, and the
reversed-index proof of Theorem 7.4]{Milnor1963}; the exhaustion
and index estimate were supplied here.

Properness of \(\overline{\mathcal M}_{0,n}\) gives its NC
boundary a stable collar. The Morse level above all critical
values also has a stable collar. These are cofinal neighbourhoods
of the same end, and their product retractions identify their
homotopy types over \(\mathcal M_{0,n}\).
Thus this boundary computes \(W_{d-1}\) and its actual homotopy
fibre. For \(d=2\) that fibre is connected; for \(d=3\) it is
simply connected. At \(n=6,k=3\) the empty multicurve gives
contractibility by Theorem~\ref{P5b1}.

The exact check \texttt{levels\_stein.py} replaces
\(|\ell|^{-4}\) by independent positive variables and expands
every nonempty principal minor over \(\Z\).
For \(n=5\) it finds five forms and 14 positive determinant
monomials; for \(n=6\), nine forms and 93 positive determinant
monomials. Every minor has constant coefficient one and positive
remaining coefficients. This verifies the displayed exhaustion
algebra; the proof, rather than a finite quotient computation,
establishes connectivity of the infinite fibres.
\end{proof}

\begin{proposition}[The higher-genus extrapolation]\label{P5b4}
\status{Refuted.}
The assertion that \(|P_k|\) is always \((k-1)\)-connected in
every genus is false. A counterexample is \((g,n,k)=(2,1,3)\).
\end{proposition}

\begin{proof}
Here \(d=4\), so \(P_3\) contains all nonempty multicurves.
Its realization is the barycentric subdivision of
\(\mathcal C_{2,1}\). By
\citet[Chapter 4, Theorem 4.1 and Corollary 4.2]{Harer1988},
this complex has the homotopy type of a nontrivial wedge of
spheres of dimension \(2g-3+n=2\).
Its nonzero second homology is the dualizing module appearing
in that corollary. Hence it is not two-connected.
The all-genera fibre identification and the contractible
endpoint \(k\geq d\) remain valid. The all-genera two-level
statement \(k\leq2\), proved next, also remains valid.
\end{proof}

\subsection{Presentations and complete first identities}

\begin{theorem}[The two-dimensional presentation procedure in all genera]
\label{P5c1}
\status{Proved.}
For every stable \((g,n)\), \(P_0\) is nonempty, \(|P_1|\) is
connected, and \(|P_2|\) is simply connected.
The action of \(\Gamma_{g,n}\) on \(|P_2|\) therefore gives a
developable complex of groups with fundamental group
\(\Gamma_{g,n}\), retaining all stabilizers.
When \(d\geq3\) the quotient uses boundary strata of dimension
at most two and their incidence flags.
Stabilizers, incidence injections, transport elements, and
triangle relations give a presentation. Dimension-two
generators can be represented by lower-dimensional tube loops;
thus dimension one supplies generators and dimension two
supplies relations. The local groups are retained with their own
presentations at every level; the level-two assertion concerns the
additional gluing relations.
\end{theorem}

\begin{proof}
For \(k\geq d\) use the empty-multicurve cone. For \(k=0\)
a pants decomposition proves nonemptiness. It remains to
treat \(1\leq k\leq2<d\), or \(k=1<d\).
The sphere dimensions of \citet[Chapter 4, Theorem 4.1]{Harer1988}
show that a curve complex of complexity one is nonempty, one of
complexity two is connected, and one of complexity at least three
is simply connected. Explicitly the sphere dimensions are \(n-4\)
for \(g=0\), \(2g-3+n\) for \(g>0,n>0\), and \(2g-2\)
for closed surfaces; each has the asserted lower bound
in the indicated stable range.

The link of a multicurve is the join of the curve complexes
of its cut components. If the residual complexities are
\(d_i\), their sum is \(d-|\sigma|\); complexity-zero
components contribute join units. When that sum is at
least two, the join is connected: either one factor
already has complexity at least two, or two nonempty
factors occur. When it is at least three, the join is
simply connected: either one factor has complexity at
least three, a connected factor joins a nonempty factor,
or at least three nonempty factors occur. These assertions
follow from the join connectivity inequality.

Now delete vertices of the barycentric curve complex in
increasing face-size order, down to sizes at least \(d-k\).
At each deletion the remaining link is the upper link,
namely the curve complex of the cut surface just considered.
Its residual complexity is at least \(k+1\), so it is
\((k-1)\)-connected. Coning it is a \(k\)-connected
attachment. The initial curve complex already has this
connectivity, and the deletion argument of
Theorem~\ref{P5b2} applies. All realizations use the weak
CW topology; relative cone attachments, or finite compact
witnesses therein, justify the infinite family of deletions.
This proves the all-genera assertions at levels zero,
one, and two.

The action on the order complex has no inversions, because
the ranks of the vertices in a chain are distinct.
Its quotient has finitely many cells, represented by
stable labelled dual graphs and contraction flags.
Take a further barycentric subdivision and choose lifts
of its cells. Let \(G_v\) be the full stabilizer of each
chosen cell barycentre. A face incidence is oriented
from the larger cell to its face, giving after transport
an injection \(\psi_a:G_{i(a)}\to G_{t(a)}\).
For composable incidences choose the comparison element
\(g_{a,b}\in G_{t(a)}\), so that
\[
                  \psi_a\psi_b
                 =\operatorname{Ad}(g_{a,b})\psi_{ab}.
\]
Use generators of all \(G_v\)'s and symbols for the
incidence arrows, with the internal stabilizer relations,
\[
 ah a^{-1}=\psi_a(h),\qquad
 ab=g_{a,b}(ab)_{\mathrm{edge}},
\]
and set a maximal tree of edge symbols equal to one.
Reversed orientations give inverse symbols.
Every \(g_{a,b}\) is determined by comparing the specified
lifts; none is silently discarded.

For completeness, a path from a chosen lift to its
translate by \(\gamma\in\Gamma_{g,n}\) gives a word in
these transport and stabilizer generators. This proves
surjectivity. A word in the kernel lifts to a closed
edge path, which bounds a simplicial disc because
\(|P_2|\) is simply connected. Removing its triangles
uses the displayed triangle relations; changing lifts
at its corners uses the stabilizer and conjugation
relations. This proves injectivity, including for
ineffective elements which act trivially on the
underlying complex but remain in its local groups.

The stabilizers are the mapping class groups of cut
surfaces with gluing twists retained, together with
the finite symmetries of their dual graphs.
The subgroup preserving all components and branches
is an extension of the product of component mapping
class groups by the gluing-twist lattice.
This extension and its coherence are retained; a
general-genus splitting is not asserted.
The component dimensions sum to at most two, so the
only positive-dimensional indecomposable factors
are of types \((0,4),(1,1),(0,5),(1,2)\).
At a dimension-two factor the already proved
connectedness of its \(P_1\) means that its boundary
tube maps onto its fundamental group. For a product
of two dimension-one factors, the join of their
nonempty boundary fibres is connected, giving the
same surjectivity. This remains true equivariantly
under graph symmetries and after taking normal-torus
bundles. Thus dimension-two generators can be replaced
by lower-dimensional words, leaving their relations.

In genus zero the dual graphs are trees, pure labels
remove the graph permutations, and framings split the
twist extensions noncanonically. The positive-dimensional
factors are then \(\mathcal M_{0,4}\),
\(\mathcal M_{0,5}\), or \(\mathcal M_{0,4}^2\).
For \(n=5\), \(P_2\) contains the empty vertex with
stabilizer \(\Gamma_{0,5}\); the abstract presentation
would be tautological there. The following independent
calculation supplies that elementary group.
\end{proof}

Write \(F_m\) for a free group of rank \(m\), and use
\(c(u,v)=v^{-1}u^{-1}vu\) throughout.
For ordinary braid generators define
\[
 A_{ij}=\sigma_{j-1}\cdots\sigma_{i+1}\,
             \sigma_i^2\,\sigma_{i+1}^{-1}
                         \cdots\sigma_{j-1}^{-1}
                  \quad (i<j).
\]
The first two finite marked points are fixed at
\(0,1\); the last spherical marked point is \(\infty\).
The remaining points move in that punctured affine line.

\begin{proposition}[An explicit presentation of \(\Gamma_{0,5}\)]
\label{P5c2}
\status{Proved.}
One has
\[
                \Gamma_{0,5}=F(x,y,z)\rtimes F(a,b),
\]
where \(a=A_{13}, b=A_{23}, x=A_{14}, y=A_{24}, z=A_{34}\).
A complete presentation consists of these five generators
and the six relations
\[
\begin{array}{lll}
 axa^{-1}=z^{-1}xz,&
 aya^{-1}=c(x,z)y c(x,z)^{-1},&
 aza^{-1}=c(x,z)z,\\[2pt]
 bxb^{-1}=x,&
 byb^{-1}=z^{-1}yz,&
 bzb^{-1}=c(y,z)z .
\end{array}
\]
The presentation complex is aspherical, with cell
counts \((1,5,6)\).
\end{proposition}

\begin{proof}
Forget the second moving point of
\(\operatorname{Conf}_2(\C\setminus\{0,1\})\).
The base is \(\C\setminus\{0,1\}\), with free
generators \(a,b\), and the fibre is
\(\C\setminus\{0,1,z_1\}\), with free generators
\(x,y,z\). This is the locally trivial fibration of
\citet[Theorem 1]{FadellNeuwirth1962}.
A section places the last point on the positive
real axis farther from the origin than every
preceding point and than \(1\). Base and fibre
are aspherical, so the homotopy exact sequence
gives the split group extension and asphericity
of its total space.

Our Artin convention is
\(\sigma_i(X_i)=X_iX_{i+1}X_i^{-1}\),
\(\sigma_i(X_{i+1})=X_i\), with other letters fixed.
For the stated kernel generators, conjugation
by a positive base braid is the Artin action of
its mirror word, obtained by negating each
Artin-generator exponent. Substituting \(A_{13}\)
and \(A_{23}\) gives the six displayed formulae.
The program \texttt{levels\_braids.py} compares
each equality using the actual braid words in
the free-group automorphism representation on
four letters. That representation is faithful;
see \citet[Theorem A, representation (1),
parameter \(1\)]{Shpilrain2001}.
The checks concern full noncommutative words.

Replace the base by a two-petalled rose and
the fibre by a three-petalled rose. Realize the
two monodromies by based graph homotopy
equivalences and form their mapping cylinders
over the base loops. The resulting graph bundle
up to homotopy has one vertex, five edges,
and exactly the six indicated two-cells.
Its universal cover is a tree of contractible
fibre trees and is contractible. The displayed
presentation is therefore complete and
aspherical.
\end{proof}

\begin{proposition}[A presentation and all first syzygies for
\(\Gamma_{0,6}\)]\label{P5c3}
\status{Proved.}
Adjoin \(p=A_{15},q=A_{25},r=A_{35},s=A_{45}\).
Then
\(\Gamma_{0,6}=F(p,q,r,s)\rtimes\Gamma_{0,5}\).
Keep the six base relators and add the twenty
relations \(huh^{-1}=\phi_h(u)\) in this table:
\[
\renewcommand{\arraystretch}{1.5}
\begin{array}{c|cccc}
h&\phi_h(p)&\phi_h(q)&\phi_h(r)&\phi_h(s)\\ \hline
a&r^{-1}pr&c(p,r)q c(p,r)^{-1}&c(p,r)r&s\\
b&p&r^{-1}qr&c(q,r)r&s\\
x&s^{-1}ps&c(p,s)q c(p,s)^{-1}
                 &c(p,s)r c(p,s)^{-1}&c(p,s)s\\
y&p&s^{-1}qs&c(q,s)r c(q,s)^{-1}&c(q,s)s\\
z&p&q&s^{-1}rs&c(r,s)s
\end{array}
\]
This gives nine generators and twenty-six
relators and a finite aspherical three-complex
with cell counts \((1,9,26,24)\).
The twenty-four identities specified below,
one for each of the six base relators and
each of \(p,q,r,s\), freely generate all
first identities modulo Peiffer identities,
equivalently the \(\pi_2\) of this presentation
complex, over \(\Z\Gamma_{0,6}\).
\end{proposition}

\begin{proof}
Forget the third moving point. The fibre is
the affine line with four points removed,
and the far-out-point section splits the
Fadell--Neuwirth extension. Substitution in
the mirrored Artin action gives the table.
Every entry is checked in
\texttt{levels\_braids.py} against actual
conjugation of \(A_{i5}\) in the five-strand
braid group. Each of the six base relators
also acts trivially on all four fibre
generators, giving twenty-four exact
free-word checks. Thus the table defines
the required action of \(\Gamma_{0,5}\);
the split extension proves completeness
of the twenty-six defining relations.

Here are the full identity words as a
finite recursive schema. Temporarily
regard all nine generators as free.
For a base word \(w\) and a fibre word
\(v\), put
\[
                  D_w(v)=wvw^{-1}\phi_w(v)^{-1}.
\]
For positive generators \(h,u\), \(D_h(u)\)
is its mixed relator. Expand other
\(D\)-words as products of conjugates
of those relators using the exact rules
\[
\begin{split}
 D_{ab}(v)&=aD_b(v)a^{-1}D_a(\phi_b(v)),\\
 D_h(vv')&=D_h(v)\phi_h(v)D_h(v')\phi_h(v)^{-1},\\
 D_h(v^{-1})&=\phi_h(v)^{-1}D_h(v)^{-1}\phi_h(v),\\
 D_{h^{-1}}(v)&=h^{-1}
                    D_h(\phi_h^{-1}(v))^{-1}h .
\end{split}
\]
The inverse automorphisms are fixed by the
mirrored-braid convention and supplied as
integer words in
\texttt{levels\_braids\_actions.json}.
Let \(R\) run through the six base words
\(hvh^{-1}(\text{its displayed image})^{-1}\),
with \(h=a,b\) and \(v=x,y,z\).
For each new generator \(u=p,q,r,s\),
\(\phi_R(u)=u\) freely, and the identity is
\[
                   R\,(uR^{-1}u^{-1})\,D_R(u)^{-1}=1. \tag{*}
\]
The first two factors are a base relator
and a conjugate inverse base relator;
the last expands into conjugates of the
twenty mixed relators by the stated rules.
This specifies every word and sign.
The program \texttt{levels\_syzygies.js} records each of
these twenty-four products and verifies
that it freely reduces to the empty word.
The products have between eight and 640
conjugate factors; the largest conjugating
word has length 171.

A check that identities hold does not
establish their completeness. For that,
let \(K_5\) be the aspherical two-complex
just constructed, and \(R_4\) a
four-petalled rose. Each component of
\[
              \Map_*(R_4,R_4)\simeq(\Omega R_4)^4
\]
is contractible, since \(R_4\) is
aspherical. Based self-homotopy
equivalences form the components indexed
by \(\operatorname{Aut}(F_4)\).
The exact action therefore constructs
a graph bundle up to homotopy over
\(K_5\), with a section. Choose the
graph maps over base edges and the
relator homotopies over base two-cells.
Contractibility of the mapping-space
components supplies their coherence.

Over a base \(q\)-cell attach its
section \(q\)-cell first. The remaining
relative pair is
\[
 (D^q\times R_4,\
   S^{q-1}\times R_4\ \cup\ D^q\times\{*\}),
\]
with quotient \(S^q\wedge R_4\), a
wedge of four \((q+1)\)-spheres.
Thus the base counts \((1,5,6)\) give
\[
 c_0=1,\quad c_1=5+4=9,\quad
 c_2=6+4\cdot5=26,\quad c_3=4\cdot6=24.
\]
The additional two-cells are the
twenty mixed relators. A base relator
disc crossed with a fibre edge has
boundary the sphere \((*)\).
The total space is aspherical by the
homotopy exact sequence of this
homotopy bundle: its base and fibre
have no higher homotopy, and its
group is the specified extension.

Let \(G=\Gamma_{0,6}\) and pass to
universal covers. There are no
cells above degree three, and
contractibility of the full cover
gives an exact sequence
\[
 0\longrightarrow \Z G^{24}
       \xrightarrow{\partial_3}\Z G^{26}
       \xrightarrow{\partial_2}\Z G^9 .
\]
Hence \(\partial_3\) is injective
onto \(\ker\partial_2\).
The universal cover of the
presentation two-skeleton is
simply connected, so Hurewicz
identifies its \(\pi_2\) with
that kernel. The twenty-four
attaching spheres are consequently
a free \(\Z G\)-basis of all
first identities modulo Peiffer
identities.
\end{proof}

This last module belongs to the
displayed finite presentation.
It is not identified with
\(\pi_2\) of the full geometric
boundary tube or of the infinite
curve complex, which have different
cell models.

\begin{theorem}[Two levels, their elementary cases, and dimension-three
identities]\label{P5c4}
\status{Proved.}
In every genus, strata of modular
dimension at most two suffice
for the ordinary fundamental-group
presentation, with full incidence
and stabilizer data.
The four indecomposable positive
dimensional types in folio 29,
scan 60, are
\[
                         (0,4),\ (1,1),\ (0,5),\ (1,2).
\]
The first two have dimension one,
the last two dimension two.

In genus zero, dimension-three
coherence has the support types
\[
 \mathcal M_{0,6},\qquad
 \mathcal M_{0,5}\times\mathcal M_{0,4},\qquad
 \mathcal M_{0,4}^3,
\]
with zero-dimensional factors allowed.
The finite presentations specified
above and below have respectively
24, 12, and 8 complete first identities.
\end{theorem}

\begin{proof}
The all-genera group presentation
is Theorem~\ref{P5c1}.
Solving \(3g-3+n=1\) in the stable
range gives \((0,4),(1,1)\);
solving \(3g-3+n=2\) gives
\((0,5),(1,2)\).
For an independent published
topological formulation,
\citet[Section 2, Theorem 2,
and Remark 4]{HatcherLochakSchneps2000}
prove simple connectivity of
the maximal-curve-system complex
with elementary moves and
relations and identify these
four support types. Their
disjoint-support commutation
relations are included.
This is an ordinary topological
two-level statement, not a
profinite or Galois reconstruction
assertion.

In genus zero each stratum is
a product of \(\mathcal M_{0,n_v}\),
with dimensions adding.
The partitions of three give
exactly the three listed types.
Theorem~\ref{P5b2} at \(k=3\)
says that this level kills
the remaining \(\pi_2\)
obstructions to its comparison.
Proposition~\ref{P5c3} gives
the first type explicitly.

For the second type use the
five-generator presentation
of \(\Gamma_{0,5}\) and add
two generators \(t_1,t_2\).
Keep its six relators and
impose the ten commutators
\([h,t_i]=1\), where
\(h=a,b,x,y,z\).
The aspherical model is
\(K_5\times R_2\), with
cell counts
\[
                              (1,7,16,12).
\]
Its twelve first identities
are \((*)\) for the six
base relators and the two
new letters, with trivial
action \(\phi_h(t_i)=t_i\).

For the third type take
generators \(a_1,a_2\),
\(b_1,b_2\), \(c_1,c_2\)
and all twelve commutators
between distinct pairs of
these three free factors.
The model \(R_2^3\) has
cell counts
\[
                               (1,6,12,8).
\]
Its eight cube identities
are the same schema for
the four base relators
\(a_i b_j a_i^{-1}b_j^{-1}\)
and each \(c_\ell\).
Both product models are
aspherical and have no cells
above dimension three.
The universal-cover and
Hurewicz argument of
Proposition~\ref{P5c3}
therefore proves that the
12 and 8 identities are
complete free bases for
their presentation modules.
The expanded-word check
verifies all of them;
the second family has
six to fourteen conjugate
factors per identity,
and the cube identities
have six each.

Normal tori and lower
pieces can already have
internal higher cells.
Thus this support statement
does not say that every
three-cell first appears at
modular dimension three.
Nor does it extend the
higher-connectivity assertion
to every genus: the
counterexample in
Proposition~\ref{P5b4}
explains the precise limit.
\end{proof}

\paragraph{Punctures and fixed boundary circles.}
The explicit sphere groups fix
punctures individually. If instead
each of the \(n\) boundary circles
is fixed pointwise, capping
introduces central boundary
twists. For \(g=0,n\geq3\)
the result is noncanonically
\(\Gamma_{0,n}\times\Z^n\).
Indeed normalization at
\(0,1,\infty\) trivializes
all individual tangent-direction
circle bundles: use \(dz\)
at finite marks and \(d(1/z)\)
at infinity. Their product
is the framed moduli space,
whose fundamental group is
the pointwise-boundary mapping
class group. Different framings
alter the splitting.

\paragraph{Computational scope.}
The supplied exact programs
verify the nonabelian hexagon
presentation, twenty-six braid
equalities, twenty-four action
compatibilities, all 44
expanded first identities for
the three dimension-three
presentations, and the Levi
principal minors.
They do not infer infinite
curve-complex connectivity
from a finite stable-tree
quotient. The connectivity
and the completeness of the
syzygies are established by
the proofs above.

\nocite{*}
\bibliographystyle{apalike}
\bibliography{EXIT_ONE_TYPE_references}
\end{document}
